% concatenated from main.tex
\documentclass[a4paper,11pt]{amsart}
\usepackage[margin=2cm]{geometry}
\usepackage[utf8]{inputenc}
\usepackage[english]{babel}
\usepackage{amssymb}
\usepackage{amsmath}
\usepackage{amsthm}
\usepackage{amsfonts}
\usepackage{array}
\usepackage{comment}
\usepackage{enumitem}
\usepackage{hyperref}
\usepackage{mathrsfs}
\usepackage{tikz}
%\usetikzlibrary{arrows,decorations.markings}
\usepackage{tikz-cd, mathtools}
%\usepackage{pxfonts}
\usepackage{mathabx}
\usepackage{xfrac}
\usepackage{graphicx} % Required for inserting images
\setcounter{MaxMatrixCols}{15}

%\usepackage[backend=biber,style=alphabetic]{biblatex}\addbibresource{Project.bib}

\theoremstyle{plain}
\newtheorem{theorem}{Theorem}[section]
\newtheorem{lemma}[theorem]{Lemma}
\newtheorem{conjecture}[theorem]{Conjecture}
\newtheorem{proposition}[theorem]{Proposition}
\newtheorem{corollary}[theorem]{Corollary}

\theoremstyle{definition}
\newtheorem{definition}[theorem]{Definition}
\newtheorem{example}[theorem]{Example}
\newtheorem{remark}[theorem]{Remark}

\newcommand{\A}{\mathbb{A}}
\newcommand{\C}{\mathbb{C}}
\newcommand{\F}{\mathbb{F}}
\newcommand{\p}{\mathbb{P}}
\newcommand{\Q}{\mathbb{Q}}
\newcommand{\R}{\mathcal{R}}
\newcommand{\Z}{\mathbb{Z}}
\newcommand{\Ox}{\mathcal{O}}

\DeclareMathOperator{\im}{im}
\DeclareMathOperator{\Pic}{Pic}
\DeclareMathOperator{\res}{res}
\DeclareMathOperator{\Spec}{Spec}
\DeclareMathOperator{\Mov}{Mov}
\DeclareMathOperator{\Proj}{Proj}

\makeatletter\let\@wraptoccontribs\wraptoccontribs\makeatother

\newcommand{\I}{\mathcal{I}}

\newcommand{\coxX}{S}
\newcommand{\pic}{\Pic}

\newcommand{\balazs}[1]{{\color{blue} [#1]}}
\newcommand{\atsushi}[1]{{\color{magenta}[#1]}}

\title{On the Cox rings of some hypersurfaces}

\begin{document}

\author{Andrew Pollock}
\address{University of Vienna, Austria}
\email{andrew.scott.pollock@univie.ac.at}

\author{Atsushi Ito}
\address{Tsukuba University, Japan}
\email{ito-atsushi@math.tsukuba.ac.jp}

\author{Bal\'{a}zs Szendr\H{o}i}
\address{University of Vienna, Austria}
\email{balazs.szendroi@univie.ac.at}

\begin{abstract}
\noindent We introduce a cohomological method to compute Cox rings of hypersurfaces in the ambient space $\p^1 \times \p^n$, which is more direct than existing methods. We prove that smooth hypersurfaces defined by regular sequences of coefficients are Mori dream spaces, generalizing a result of Ottem. We also compute Cox rings of certain specialized examples. In particular, we compute Cox rings in the well-studied family of Calabi--Yau threefolds of bidegree $(2,4)$ in $\p^1 \times \p^3$, determining explicitly how the Cox ring can jump discontinuously in a smooth family. 
\end{abstract}

\maketitle

\thispagestyle{empty}

\section*{Introduction}

Let $X$ be a smooth projective variety over~$\C$, with Picard group $\Pic X$ a free $\Z$-module for simplicity. The Cox ring 
\begin{equation*}
    \R(X) = \bigoplus_{D \in \Pic X} H^0(X,\Ox_X(D))
\end{equation*}
has the natural structure of a $\C$-algebra graded by $\Pic X$; details including a careful definition can be found in \cite{CoxBook}. Let $\iota\colon Z\hookrightarrow X$ be a smooth or mildly singular hypersurface, defined by a homogeneous equation $r\in 
\R(X)$. Under suitable conditions, the restriction map $\Pic(X)\to\Pic(Z)$
is an isomorphism; we will assume this for the rest of the Introduction, and identify the Picard groups via restriction. 
Alongside the restriction map on line bundles, we also have a restriction map between Cox rings
\[\iota^*\colon \R(X)\to \R(Z). 
\]
This map was first studied explicitly by Hausen and Artebani--Laface~\cite{hausen, AL}, giving conditions for $\iota^*$ to be surjective, leading to an isomorphism
\[\R(Z) \cong \R(X)/\langle r \rangle.\]
The aim of our paper is to give a new, explicit method to compute the Cox rings of some hypersurfaces~$Z$ in the specific ambient space $X = \p^1 \times \p^n$ for $n\geq 3$. In this case, the map $\iota^*\colon \R(X)\to \R(Z)$ is not surjective, with $O_Z(D|_Z)$ having sections that do not come from $O_X(D)$ for certain $D$.

We first define the hypersurfaces studied in this paper. Throughout, we fix $n \geq 3$ and $d,e \geq 1$.

\begin{definition}\label{def:Ztypes}
Let $Z$ be an irreducible and normal hypersurface
\begin{equation} \label{eq:generalZform}
Z = \left\{g_0x_0^d + g_1 x_0^{d-1} x_1 + \ldots + g_d x_1^d=0\right\}\subset\p^1\times\p^n
\end{equation}
defined by a set of homogeneous degree $e$ polynomials $g_0,\dots,g_d \in \C[y_0,\dots,y_n]$.
    \begin{enumerate}
        \item[(i)] We say $Z$ admits an {\em index sequence of length $l$}, with $1 \leq l \leq d$, if there are indices $0 = i_0 < \cdots < i_l = d$ such that $g_i \neq 0$ if and only if $i \in \{i_0,\dots,i_l\}$; we call the set of indices $\{i_0,\dots,i_l\}$ the {\em index sequence} of $Z$.
        \item[(ii)] We call $Z$ {\em regular with index sequence $\{i_0,\ldots, i_l\}$}, if it is non-singular, admits the index sequence $\{i_0,\dots,i_l\}$, and the homogeneous polynomials $g_{i_0},\dots,g_{i_l}$ form a regular sequence in the graded ring $\C[y_0,\dots,y_n]$.
    \end{enumerate}
\end{definition}

%\begin{definition}
%    Let $Z \subset \p^1 \times \p^n$, as in \eqref{eq:generalZform}. Suppose there is a partition $0 = i_0 < \dots < i_l = d$ for some $0 \leq l \leq d$, such that $g_i \neq 0$ if and only if $i \in \{g_0,\dots,g_d\}$ and $g_{i_0},\dots,g_{i_l}$ forms a regular sequence in $\C[y_0,\dots,y_n]$. We say that $Z$ admits a partition in a regular way and call $(i_k)_{k=0}^{l}$ the partition (of $d$) associated to $Z$. 
%\end{definition}

Note that 
% if $Z$ is of regular general or special type then it admits a partition in a regular way. Furthermore, 
amongst hypersurfaces $Z\subset\p^1\times\p^n$ admitting a given index sequence $\{i_0,\ldots, i_l\}$, being regular with index sequence $\{i_0,\ldots, i_l\}$ is an open condition. 

The following is our main result (see Theorem~\ref{thm:general}, Proposition~\ref{prop:converse}). 

\begin{theorem} \label{thm:mainintro}
Suppose that $Z\subset\p^1\times\p^n$ is a nonsingular hypersurface of the form \eqref{eq:generalZform}, 
with $n \geq 3$ and $d,e \geq 1$, admitting an index sequence $\{i_0,\dots,i_l\}$ of length $l$.
\begin{enumerate}
\item[(i)] Suppose that $Z$ is regular with index sequence $\{i_0,\dots,i_l\}$, that is $g_{i_0},\dots,g_{i_l}$ form a regular sequence in $\C[y_0,\dots,y_n]$. Then there is an isomorphism of $\Z^2$-graded algebras
\begin{equation}\label{presentationR}
    \R(Z)\cong \C[x_0,x_1,y_0,\dots,y_n,w_1,\dots,w_l]/I,
    \end{equation}
    where $\deg(x_i) = (1,0)$, $\deg (y_j) = (0,1)$, $\deg(w_k) = (i_{k-1}-i_k,e)$ for relevant $i,j,k$ and $I$ is the homogeneous ideal
    \[
    I = \langle x_1^{i_1-i_0} w_1 + g_0, x_1^{i_2-i_1} w_2 + g_1 - x_0^{i_1-i_0} w_1,\ldots, g_d - x_0^{i_l-i_{l-1}} w_l\rangle.
    \]
    In particular, $Z$ is a Mori dream space. 
\item[(ii)] Conversely, suppose that there exists an isomorphism of the form~\eqref{presentationR}, with degrees and ideal $I$ as given. Then $g_{i_0},\dots,g_{i_l}$ form a regular sequence in $\C[y_0,\dots,y_n]$. In other words, $Z$ is regular with index sequence $\{i_0,\dots,i_l\}$.
\end{enumerate}
\end{theorem}

This result substantially strengthens a result of Ottem~\cite{ottem}, who proves (i) only for general $Z$ admitting full index sequence $\{0,1,\dots,d\}$. It also provides a converse statement.
%Note that $d=0$ is special here, in this case $Z=\{g_0=0\}\subset\p^1\times\p^n$ is itself a product, there are no new $z_i$ variables, and we have an isomorphism $\R(Z) \cong \R(\p^1\times\p^n)/\langle g_0 \rangle$; we only included this case in the statement above as it provides a convenient base for our induction.  

\begin{comment}
 The following is our main result (see Theorem~\ref{thm:general}).

\begin{theorem}\label{thm:mainintro}
    Fix $n \geq 4$ and $d \geq 0,e \geq 1$ or $n = 3$ and $d,e \geq 1$. Suppose that $Z \subset \p^1 \times \p^n$ is either of regular general type or special type and let $(i_k)_{k=0}^{l}$ be its associated partition. Then we have an isomorphism of $\Z^2$-graded algebras
\begin{equation}\label{presentationRwithw}\R(Z)\cong \C[x_0,x_1,y_0,\dots,y_n,w_1,\dots,w_l]/I,\end{equation}
where $\deg(x_i) = (1,0)$, $\deg (y_j) = (0,1)$, $\deg(w_k) = (i_{k-1}-i_k,e)$ for relevant $i,j,k$ and $I$ is the bihomogeneous ideal
\[
\langle x_1^{i_1 - i_0} w_1 + g_{i_0}, x_1^{i_2 - i_1} w_2 + g_{i_1} - x_0^{i_1 - i_0} w_1,\ldots, g_{i_l} - x_0^{i_l - i_{l-1}} w_l\rangle.
\]
In particular, $Z$ is a Mori Dream Space. 
\end{theorem}

For the case where $Z$ is of regular general type, Theorem~\ref{thm:mainintro} is a strengthening of a result of Ottem \cite{ottem}, who makes a genericity assumption. We show that regularity of the coefficients $g_0,\dots,g_d$ in $\C[y_0,\dots,y_n]$ is sufficient for the Cox ring to take the form \eqref{presentationRwithw}. We are able to complete this further with the following converse result (see Proposition~\ref{prop:conversemain}).

\begin{theorem}\label{thm:introconverse}
    Suppose that $Z$ is any hypersurface of the form \eqref{eq:generalZform}. Suppose there is an isomorphism of graded algebras of the form
\begin{equation}\label{presentationRwithz}
    \R(Z)\cong \C[x_0,x_1,y_0,\dots,y_n,z_1,\dots,z_d]/\langle x_1 z_1 + g_0, x_1 z_2 + g_1 - x_0 z_1,\ldots, g_d - x_0 z_d\rangle,
    \end{equation}
    where $\deg(x_i) = (1,0)$, $\deg (y_j) = (0,1)$, $\deg(z_k) = (-1,e)$ for relevant $i,j,k$. Then $g_0,\dots,g_d$ forms a regular sequence in $\C[y_0,\dots,y_n]$. In particular, if $Z$ is non-singular then it is of regular general type.
\end{theorem}

The first two authors proved Theorem~\ref{thm:mainintro} for hypersurfaces $Z$ of regular type in a first version of this paper. Thanks to work of Atsushi Ito, we have extended the result to $Z$ of special type $Z$. We expect the result to hold only under the assumption that $Z$ is irreducible, non-singular and admits a partition in a regular way, but are as yet unable to prove this. If $l \neq d$, our proof requires the nonsingularity of the varieties $Z(k)$, as defined in Definition~\ref{def:Ztypes}(ii), to show that we have found all relations in the Cox ring. Our proof thus suffices to show that an irreducible, non-singular $Z$ admitting a partition in a regular way has Cox ring a quotient of the presentation in \eqref{presentationRwithw}, and in particular $Z$ is a Mori Dream Space. The appearace of the varieties $Z(k)$ foreshadows a generalisation of Theorem~\ref{thm:mainintro} to hypersurfaces of a $\p^1 \times W$, with $W$ a complete intersection. This is will feature in a follow up paper by the first author.
\end{comment} 

Our proofs use standard cohomological machinery to reduce the problem to an algebraic one, involving syzygies of the set $\{g_{i_0}, \dots, g_{i_l}\}$ of elements of $R$, as well as induction. Our arguments are longer than those of~\cite{ottem}, but are more direct, able to handle also the more general setting of Theorem \ref{thm:mainintro}.  More details, and further results based on the method, will be presented in~\cite{pollock}. Our examples can also be studied using work of Herrera, Laface and Ugaglia~\cite{lafaceetal}, who use yet another method involving localizations of the Cox ring $\R(X)$ of the embedding space $X=\p^1\times\p^n$; our approach gives a natural reason for the appearance of these localizations. 

For $n=3$ and $(d,e)=(2,4)$, the family 
\[{\mathcal Z} = \left\{g_0 x_0^2 + t g_1 x_0x_1 + g_2 x_1^2 =0 \right\}\subset \p^1\times \p^3\times \A^1_t
\]
is a well-studied family of Calabi--Yau threefolds, see in particular the recent~\cite{constantinetal, constantin}. In this case, Theorem~\ref{thm:mainintro} shows that this is an example of a smooth family of Calabi--Yau threefolds in which the Cox ring jumps on a closed subset of the moduli space; see Theorem~\ref{thm_cy}. We also discuss a further, singular example in this family in Example~\ref{ex:det} .

Theorem~\ref{thm:mainintro} applies in particular when $Z$ is regular and $e=1$,
with the projection $Z\to\p^1$ being a $\p^{n-1}$-bundle. If $d>n$, then the linear forms $\{g_0, \dots, g_d\}$ 
are even linearly dependent, and certainly cannot form a regular sequence, so by Theorem~\ref{thm:mainintro}(ii), the Cox ring is not of the form~\eqref{presentationR}, contrary to the statement \cite[Thm.1.1(iv)]{ottem}. We discuss one case in Example~\ref{ex:linear}-Proposition~\ref{prop:linear}.

Finally in Example~\ref{ex:regcasemodel}, we explain from our biregular point of view a birational transformation $Z\dashrightarrow Z'$ that played an important role in Ottem's arguments~\cite[Section 2]{ottem}.

Section 1 explains our cohomological approach to the problem. Section 2 is the main part of our paper, where we build up to the proof of our main results, modulo an algebraic statement that is relegated to Section 4. Section 3 discusses our examples.

\medskip

\noindent {\bf Acknowledgments and authorship } \ We thank Antonio Laface for several comments and for an argument sketched in Section 3, %Atsushi Ito for comments on a first version of the paper that lead to this second version, 
and Geoffrey Mboya for discussions about Cox rings. This paper forms part of the doctoral project of the first-named author, under the supervision of the last-named author. After an initial preprint by these two authors, the middle author made suggestions for strengthening the results. A proof of this stronger Theorem~\ref{thm:mainintro} was worked out jointly, leading to the current manuscript with an updated set of authors. 

\section{Our approach}
\subsection{Basics}

Recall our setup from the Introduction: assume that $\iota\colon Z\hookrightarrow X$ is a smooth or mildly singular hypersurface, with an isomorphism $\Pic(X)\cong\Pic(Z)$, a finitely generated free abelian group that we will denote by $\pic$. 
Let $E=[Z]\in \pic$ be the class of~$Z$. 

The Cox rings $\R(X)$ and $\R(Z)$ are both $\pic$-graded algebras
\[\R(X) = \bigoplus_{D\in\Pic} \R(X)_D, \ \ \ \R(Z) = \bigoplus_{D\in\Pic} \R(Z)_D, \]
connected by the graded homomorphism $\iota^*\colon \R(X)\rightarrow \R(Z)$. 
The interesting case for us is when the inclusion $\iota^*\R(X)\hookrightarrow \R(Z)$ is not surjective. 

The hypersurface $Z$ is given by a section $r \in \R(X)_E$. For $D\in \pic$, consider the standard exact sequence
\begin{equation}\label{seq:standard}
0\longrightarrow \Ox_X(D-E)\stackrel{\cdot r}\longrightarrow \Ox_X(D) \stackrel{\iota^*}\longrightarrow \Ox_Z(D)\longrightarrow 0. \end{equation}
The associated long exact sequence yields the short exact sequence of vector spaces
\begin{equation}\label{seq:short}
0\longrightarrow \iota^*\R(X)_D\stackrel{}\longrightarrow \R(Z)_D \stackrel{\delta}\longrightarrow N_D\longrightarrow 0,
\end{equation}
where $\iota^*\R(X)_D\cong \R(X)_D / r \cdot \R(X)_{D-E}$, and $N_D=\ker r_{D*}$ is the kernel of the induced map on first sheaf cohomology
\begin{equation}\label{seq:r*}
r_{D*} : H^1(X, \Ox_X(D-E)) \longrightarrow H^1(X, \Ox_X(D)).
\end{equation}
It will sometimes be natural to put all these maps together, to get the map
\begin{equation}
r_{*} : \bigoplus_{D\in\Pic} H^1(X, \Ox_X(D-E)) \longrightarrow \bigoplus_{D\in\Pic} H^1(X, \Ox_X(D)).
\end{equation}
We have $\iota^*\R(X)= \R(Z)$ if and only if the map $r_*$ is injective. 

\subsection{\v Cech cohomology considerations}\label{subsec:cech}

Fix a totally ordered affine cover $\mathcal{U} = \{U_i\}_{i \in I}$ of~$X$. 
From~\eqref{seq:standard}, we get the double complex

\begin{equation}\label{complex}\begin{tikzcd}
{} & 0\arrow{d} & 0\arrow{d} \\
0\arrow{r} & \displaystyle\prod_{i}\Ox_X(D-E)(U_i) \arrow{r}\arrow{d}{\cdot r} &  \displaystyle\prod_{i<j}\Ox_X(D-E)(U_i \cap U_j)\arrow{r} \arrow{d}{\cdot r}& \ldots \\
0\arrow{r} & \displaystyle\prod_{i}\Ox_X(D)(U_i)\arrow{r}{d_0} \arrow{d}{\iota^*}& \displaystyle\prod_{i<j}\Ox_X(D)(U_i \cap U_j)\arrow{r} \arrow{d}{\iota^*} & \ldots \\
0\arrow{r} &\displaystyle\prod_{i}\Ox_Z(D)(U_i)\arrow{r} \arrow{d}& \displaystyle\prod_{i<j}\Ox_Z(D)(U_i \cap U_j)\arrow{r} \arrow{d} & \ldots \\
{} & 0 & 0 
\end{tikzcd}\end{equation}

\begin{lemma} Consider \v Cech $0$-cochains $(v_i) \in \prod_{i}\Ox_X(D)(U_i)$. Call a $0$-cochain $(v_i)$ compatible (with respect to $Z$) if for each $i < j$ there exists $q_{i,j} \in \Ox_X(D-E)(U_i \cap U_j)$ such that $v_i|_{U_i \cap U_j} - v_j|_{U_i \cap U_j} = q_{i,j}r$. Call compatible $0$-cochains $(v_i), (v'_i)$ equivalent, if for each $i$, there exists $q_i \in \Ox_X(D-E)(U_i)$ such that $v'_i - v_i = q_i r$. 
\begin{enumerate}
    \item[(i)] Elements $v\in \R(Z)_D$ are naturally in bijection with equivalence classes of compatible $0$-cochains
    $$(v_i)\in\prod_{i}\Ox_X(D)(U_i).$$ 
    Under this equivalence, the restriction of a section $g \in \R(X)_D$ corresponds to the \v Cech $0$-cocycle $(g|_{U_i}) \in \prod_{i}\Ox_X(D)(U_i)$. 
    \item[(ii)] Multiplication in the algebra $\R(Z)$ is compatible with multiplication of cochains: if $D,D' \in \pic$ and $v \in \R(Z)_D$, $v' \in \R(Z)_{D'}$ are represented by compatible $0$-cochains $(v_i), (v'_i)$ respectively, then $v v' \in \R(Z)_{D+D'}$ is represented by the compatible $0$-cochain $(v_i v'_i)$.
    \end{enumerate}\label{lemma_cocycle}
\end{lemma}
\begin{proof} These statements follow from considering the first column of~\eqref{complex}, recalling that the global section functor on sheaves is exact over affine schemes. Note that a $0$-cochain $(v_i)$ restricts to a $0$-cocycle on $Z$ if and only if it is compatible.
\end{proof}

We are interested in cases when there are equivalence classes of compatible $0$-cochains $(v_i)$ that do not come from the restriction of a section $g\in \R(X)$, so that $(v_i)$ is not a $0$-cocycle.
For $D\in\Pic$, in order to find a complementary vector space to $\iota^*\R(X)_D$ in $\R(Z)_D$, we need to complete the first three steps below.
\begin{enumerate}
    \item Give an explicit description of the first sheaf cohomologies of $\Ox_X(D-E)$ and $\Ox_X(D)$ and of the map $r_{D*}$ of~\eqref{seq:r*}. 
    \item Describe the kernel $N_D=\ker r_{D*}$ as a subspace of $H^1(X, \Ox_X(D-E))$.
    \item Choose a right inverse $\epsilon\colon N_D\to \R(Z)_D$ of $\delta$ that lifts elements of $N_D$ to sections on~$Z$.
    \item Study the ring structure of the direct sum of the resulting spaces to reconstruct the whole structure of~$\R(Z)$. 
\end{enumerate}
Steps (1) and (2) have to be completed in particular cases of interest by explicit calculation. Finding a map $\epsilon$ in step (3) comes down to untangling the connecting homomorphism $\delta$. Given $s \in N_D$ for some $D \in \pic$, let 
\[(s_{i,j})\in \prod_{i<j}\Ox_X(D-E)(U_i \cap U_j)\] be a representative $1$-cocycle. Then $(r \cdot s_{i,j})$ is cohomologically trivial, thus the image of a $0$-cochain \[(v_i) \in \prod_{i}\Ox_X(D)(U_i),\] which is compatible since $v_i|_{U_i \cap U_j} - v_j|_{U_i \cap U_j} = s_{i,j}r$ for each $i < j$. This compatible $0$-cochain corresponds via Lemma~\ref{lemma_cocycle} to an element $v \in \R(Z)_D$ that satisfies $\delta v = s$. Note that the element $v$ is only well-defined up to an element in $\iota^* \R(X)_D$, since $(v_i)$ is well-defined only up to an element of $\ker d^0 = \R(X)_D$. To construct a map $\epsilon$, we need to make compatible choices so that we indeed end up with a well-defined lifting map
\[
\epsilon \colon  N_D  \to  \R(Z)_D,
\]
a right inverse to $\delta$. To complete step (4), further arguments are needed that will be based on Lemma~\ref{lemma_cocycle}(ii). 

To conclude the general discussion, let us make one further remark. Suppose that $v \in \R(Z)_D$ is represented by a compatible $0$-cochain $(v_i)$ via Lemma~\ref{lemma_cocycle}(i). For fixed $i \in I$, we have $v_i \in \Ox_{X}(U_i)(D)$, and thus $v_i = \frac{h_i}{g_i}$ for some regular functions $h_i,g_i$ on $X$ with $g_i$ not vanishing on $U_i$. Identifying $h_i$ and $g_i$ with their restrictions on $Z$, the section $g_i v_i - h_i$ is zero on $Z \cap U_i$, which is open and dense in $Z$. For each $i$, we thus obtain an equation
\[g_i v - h_i = 0\]
in $\R(Z)$. This explains the relevance of localisations of~$\R(X)$ for the problem, which is the basis of the approach taken in~\cite{lafaceetal}. 

\section{Cox rings of hypersurfaces in a product of projective spaces}
\subsection{Basics}
Fix $n \geqslant 3$ and let $X = \p^1 \times \p^n$. We use homogeneous coordinates $x_0, x_1$ and $y_0, \dots, y_n$ on $\p^1$ and $\p^n$ respectively. If $p_1, p_2$ are the projection maps to the respective factors $\p^1, \p^n$, then $\Pic X\cong \Z^2$ is generated by the pullbacks $p_1^* \Ox_{\p^1}(1)$ and $p_2^* \Ox_{\p^n}(1)$ respectively. The Cox ring of $X$ is \[\coxX=\R(X) = \C[x_0,x_1,y_0,\dots,y_n],\] with the $\Z^2$-grading given by $\deg(x_i) = (1,0)$ and $\deg(y_j) = (0,1)$. The ring \[R=\R(\p^n) = \C[y_0, \dots, y_n]\] is naturally a $\Z^2$-graded subring of~$S$, making~$S$ into a $\Z^2$-graded $R$-module.

Fix 
%$n \geq 4$ and $d \geq 0,e \geq 1$ or 
$n \geq 3$ and $d,e \geq 1$. Let $r \in \coxX_{(d,e)}$ and
\[Z=\left\{ r= 0\right\} \subset X\] 
be the corresponding hypersurface of bidegree $(d,e)$ in $X$. Then $Z$ is defined by the bihomogeneous section 
\begin{equation}
    r = g_0 x_0^d + g_1 x_0^{d-1} x_1 + \ldots + g_d x_1^d
\end{equation}
for some $g_0, g_1, \ldots, g_d \in R_e$. If $Z$ is nonsingular, the Lefschetz hyperplane theorem shows that the inclusion $\iota\colon Z\hookrightarrow X$ induces an isomorphism of Picard groups $\iota^*\colon\Pic X \cong \Pic Z \cong\Z^2$ that will be denoted $\Pic$. We will also consider a singular case, where we will comment on the relation between the Picard groups explicitly.

We require an ordered open cover $\mathcal{W}$ of $X$. For $i = 0,1$, let $U_i \subset \p^1$ be the standard affine open defined by $x_i \neq 0$. For $j = 0,\dots,n$, let $V_i \subset \p^n$ be the standard affine open defined by $y_j \neq 0$. Set $W_{ij} = U_i \times V_j$ for each $i,j$ and let $\mathcal{W} = \{W_{ij}\}_{i,j}$, with the lexicographic ordering. If $ij < kl$, then denote $W_{ij} \cap W_{kl}$ by $W_{ij,kl}$.

\subsection{The degeneracy locus of the second projection}
%Suppose that $d \geq 1$. 
Consider the projection $p_2\colon X\to\p^n$ restricted to $Z$. This is generically a $d$-to-one cover, but it has a degeneracy locus 
\[ Y = \{ g_0=g_1=\cdots = g_d=0\}\subset\p^n,\]
over which the fibres are isomorphic to~$\p^1$. Let 
\[ I_Y = \langle g_0, g_1, \ldots g_d\rangle \lhd R = \C[y_0, \ldots, y_n]\]
be the corresponding ideal. The quotient $R/I_Y$ has a finite free $R$-module resolution of the form 
\begin{equation} \label{eq:sequence_for_T} \ldots \longrightarrow R^{\gamma}  \stackrel{C}\longrightarrow R^{\beta}\stackrel{B}\longrightarrow R^{d+1}\stackrel{A}\longrightarrow R \longrightarrow R/I_Y\longrightarrow 0,\end{equation}
where we suppressed degrees, with $A=(g_0, g_1,\ldots , g_d)$, and $B,C$ matrices of homogeneous elements of the appropriate sizes and degrees. 

Of particular interest to us is the case where $Z$ is regular with index sequence $\{i_0,\ldots, i_l\}$. It is easy to check that there are always such choices with $Z \subset X$ nonsingular. Then \eqref{eq:sequence_for_T} is the sum of the standard Koszul complex for $\{g_{i_0},\dots,g_{i_l}\}$ and a trivial complex. The locus $Y \subset \p^n$ is a complete intersection of codimension $l+1$, and the inverse image $E = p_2^{-1}(Y) \subset X$ has codimension $l$ in $Z$. In the extreme case, where $l=1$ and $r = g_0 x_0^d + g_d x_1^d$, the variety $E$ is a divisor ruled over $Y$. 

\subsection{The case $d=0$}
The case $d=0$ is somewhat degenerate, but it will be useful to include as it provides a convenient case for our induction. For $d=0$, $Z \cong \p^1 \times Z_1$, with $Z_1= \{g_0=0\} \subset \p^n$. If moreover $n \geq 4$, then clearly $\R(Z) \cong \R(\p^1\times\p^n)/\langle g_0\rangle$, compatibly with the statement of Theorem~\ref{thm:mainintro}. However, if $(n,d) = (3,0)$, then $\Pic X$ and $\Pic Z$ may not be isomorphic and Theorem~\ref{thm:mainintro} may fail to hold. Indeed, for example if $(n,d,e) = (3,0,2)$ and $g_0$ is a general quadric, then \[Z\cong \p^1\times Q \cong \p^1 \times \p^1 \times \p^1\]
with $Q\subset\p^3$ a general quadric surface, so that there is a $\Z^3$-graded isomorphism
    \[
    \R(Z) \cong \C[x_0,x_1,u_0,u_1,v_0,v_1].
    \]
However, the Veronese subring of $\R(Z)$ associated to the subgroup $\iota^* \R(X) \subset \R(Z)$, which is generated by $(1,0,0)$ and $(0,1,1)$, takes the form \eqref{presentationR}, which in this case is
 \[  \bigoplus_{D\in\iota^*\!\Pic X}\R(Z)_D \cong k[x_0,x_1,y_0,y_1,y_2,y_3]/(g_0). \]

\subsection{Describing the maps on first cohomology}

The first \v Cech cohomology groups of $X=\p^1\times\p^n$ can be written using the Künneth formula, or computed explicitly using the cover $\mathcal{W}$. This gives the following result, completing Step (1) of 
our plan sketched in~\ref{subsec:cech}. 
\begin{lemma} \label{lem:H1}\begin{enumerate} \item[(i)] We have natural $\Z^2$-graded vector space isomorphisms 
\begin{equation}\label{eq:isoH1}
    \bigoplus_{(a,b) \in \Z^2} H^1(X,\Ox_X(a,b)) \cong \frac{1}{x_0 x_1}R[x_0^{-1}, x_1^{-1}] \cong S[x_0^{-1}, x_1^{-1}]/L,
\end{equation}
where recall $R=\C[y_0,\ldots, y_n]$, $S=\C[x_0,x_1, y_0,\ldots, y_n]$, and $L$ is the vector subspace of 
$S[x_0^{-1}, x_1^{-1}]$ generated by all monomials in $x_i, y_j$ in which $x_0$ or $x_1$ appear with non-negative exponent.
\item[(ii)] The map
\[
r_{*} : \displaystyle\bigoplus_{D\in\Pic} H^1(X, \Ox_X(D-E)) \longrightarrow \displaystyle\bigoplus_{D\in\Pic} H^1(X, \Ox_X(D))
\]
can be identified with the map 
\[\begin{array}{rcl}
r_{*} : S[x_0^{-1}, x_1^{-1}]/L & \longrightarrow & S[x_0^{-1}, x_1^{-1}]/L\\
f & \mapsto &  r \cdot f \mod L
\end{array}
\]
that multiplies an element $f\in S[x_0^{-1}, x_1^{-1}]/L\cong\frac{1}{x_0 x_1}R[x_0^{-1}, x_1^{-1}]$ by $r$, and ignores those terms in which $x_0$ or $x_1$ occur with non-negative exponent. 
\end{enumerate}
\end{lemma}
\begin{proof} This first isomorphism in~\eqref{eq:isoH1} is clear from the K\" unneth formula. This identifies an element $f \in \frac{1}{x_0 x_1}R[x_0^{-1}, x_1^{-1}]_{(a,b)}$ with class of the $1$-cocycle $(s_{ij,kl})$ representing an element of $H^1(X,\Ox_X(a,b))$, where $s_{ij,kl} = f$ if $i \neq k$ and $s_{ij,kl} = 0$ if $i = k$. The second isomorphism in~\eqref{eq:isoH1} is immediate. The description in (ii) is also clear from the explicit cocycle description. 
\end{proof}
Note that the spaces appearing in the isomorphism~\eqref{eq:isoH1} naturally have the structure of a graded $R$-module, which we will use in our subsequent arguments. 

We move on to Step (2) of our plan. Since $H^1(X,\Ox_X(a,b))$ is only non-zero for $a<0$, we will use the index $-a$ rather than $a$. For fixed $a > 0$, there is an isomorphism of $R$-modules
\begin{equation}\label{eq:isodirectsum} \begin{array}{rcl}R^{a-1}& \stackrel{\sim}\longrightarrow& \displaystyle\bigoplus_{b \in \Z} 
\frac{1}{x_0 x_1}R[x_0^{-1},x_1^{-1}]_{(-a,b)} \\
(f_1, \ldots, f_{a-1}) & \mapsto & f=\displaystyle\sum_{k=1}^{a-1} \frac{f_k}{x_0^{a-k} x_1^k}.
\end{array}\end{equation}
Using Lemma~\ref{lem:H1}(ii), we obtain the following description. 
\begin{proposition}\label{prop:ker} For $D=(-a,b)\in\Pic$, the kernel $N_{(-a,b)}=\ker r_{D*}$ can only be nonzero if $-a \leq d-2$. 
\begin{enumerate}\item[(i)]
For~$-1 \leq -a \leq d-2$, we have a graded isomorphism
\begin{equation*}
    \bigoplus_{b \in \Z} N_{(-a,b)} \cong R^{a+d-1}[e], 
\end{equation*}
where $[e]$ denotes a shift in grading so that, for example, 
$N_{(d-2,b)} \cong R_{b-e} \cdot \frac{1}{x_0 x_1}$
for $b \in \Z$. 
\item[(ii)]
For~$-a<-1$, we have a graded isomorphism
\begin{equation*}
    \bigoplus_{b \in \Z} N_{(-a,b)} \cong \ker A_a[e],
\end{equation*}
where $A_a$ is the $R$-module homomorphism
\[
\begin{array}{rcccl}
  A_a & \colon & R^{a+d-1} & \longrightarrow & R^{a-1} 
  \\  && (f_k)_{k=1}^{a+d-1} & \longmapsto &  (g_0 f_{k} + g_1 f_{k+1} + \ldots + g_d f_{k+d})_{k=1}^{a-1}
\end{array}
\]
defined by the matrix
\begin{equation*}
         \begin{pmatrix}
            g_0 & g_1 & g_2 & \cdots & g_d & 0 & 0 & \cdots \cdots & 0 & 0 & 0 & \cdots & 0 & 0\\
            0 & g_0 & g_1 & \cdots & g_{d-1} & g_d & 0 & \cdots \cdots & 0 & 0 & 0 & \cdots & 0 & 0\\
            \ & \vdots & \ & \ddots & \ & \vdots & \ & \ddots & \vdots & \ & \ & \ddots & \ & \vdots\\
            0 & 0 & 0 & \cdots & 0 & 0 & 0 & \cdots \cdots & 0 & g_0 & g_1 & \cdots & g_d & 0 \\
            0 & 0 & 0 & \cdots & 0 & 0 & 0 & \cdots \cdots & 0 & 0 & g_0 & \cdots & g_{d-1} & g_d
        \end{pmatrix}.
    \end{equation*}
\end{enumerate}
% For $a>1$, it is given by the $(a-1)\times (a+d-1)$ matrix $(a_{st})$ with entries defined by
% \begin{equation*}
%a_{st} =
% \left\{ \begin{array}{ll}	 g_l & \mbox{if } 0 \leq l \leq d, t-s = l  \\
%			 0 & \mbox{otherwise.}  \\
%\end{array}		\right.\end{equation*}
\end{proposition}

The first non-trivial example of a map $A_a$ is
\begin{equation*}\begin{array}{rcccl}
    A_2 & \colon & R^{d+1} & \longrightarrow & R \\
     && (f_1, \ldots, f_{d+1}) & \mapsto & g_0 f_{1} + g_1 f_{2} + \ldots + g_d f_{d+1}
\end{array}     
\end{equation*}
which is the same as the map $A$ in the complex~\eqref{eq:sequence_for_T}, and is the first map in the Koszul complex associated with the sequence $\{g_0, \dots, g_d\}$ in $R$. Its kernel, the {\em syzygies} between the defining polynomials $\{g_0, \dots, g_d\}$, give us elements in $\R(Z)_{(-2,b)}$ for different values $b$. %In some cases, generators thus obtained (once lifted back to the Cox ring as explained in the next section) will be sufficient to generate the Cox ring $\R(Z)$. In particular, a distinguished case will be when $\{g_0, \dots, g_d\}$ form a regular sequence in $R$. 

\subsection{Lifting cohomology elements to sections in the Cox ring} \label{sec:lifting} In this subsection, we describe how to perform Step~(3) from Section~\ref{subsec:cech} in our context. For $D=(-a,b)\in\Pic$, recall the standard exact sequence~\eqref{seq:short}
\[
0\longrightarrow \iota^*\R(X)_{(-a,b)}\stackrel{}\longrightarrow \R(Z)_{(-a,b)} \stackrel{\delta}\longrightarrow N_{(-a,b)}\longrightarrow 0.
\]
We need to construct a lifting map $\epsilon\colon N_{(-a,b)} \rightarrow \R(Z)_{(-a,b)}$ that is a right inverse to $\delta$. By Proposition~\ref{prop:ker}, $N_{(-a,b)}$ can only be non-zero if $-a \leq d-2$; fix such an $a$. Let $f \in N_{(-a,b)}$. Then $u = r \cdot f$ is a degree $(-a,b)$ element in $S[x_0^{-1}, x_1^{-1}]$ such that in each non-zero term, at least one of $x_0, x_1$ appears with non-negative (possibly zero) exponent. Let $u^{(0)}$ be the sum of terms in $u$ where $x_0$ has negative exponent, and let $u^{(1)}$ be the sum of terms in which $x_1$ has negative exponent. Let $u^{(2)}$ be the sum of remaining terms of $u$. Then 
\begin{equation} \label{eq:split}
u  = u^{(0)}+u^{(1)}+u^{(2)}
\end{equation}
with
\[ \begin{array}{c}
    u^{(i)} \in \Ox_X(-a,b)(W_{ij})  \mbox{ for all } j \mbox{ and } i = 0,1, \\
    u^{(2)} \in S_{(-a,b)}.
\end{array}
\]
Then the $0$-cochain $(v_{ij})$ on $X$ given by 
\begin{equation}\label{eq_nonuniquelift}
v_{ij} =
 \left\{
  \begin{array}{ll}
			 u^{(0)} + \frac{1}{2} u^{(2)} & \mbox{if } i = 0, \\
			 -u^{(1)} - \frac{1}{2} u^{(2)} & \mbox{if } i = 1, 
		\end{array}
		\right.
\end{equation}
is compatible, thus by Lemma~\ref{lemma_cocycle} gives a section $v \in \R(Z)_{(-a,b)}$. The considerations at the end of Section~\ref{subsec:cech} translate into 
\begin{proposition} \label{prop:epsilon}
    The map 
\[\begin{array}{rcccl} \epsilon& \colon &  N_{(-a,b)} & \to & \R(Z)_{(-a,b)}\\
&& f & \mapsto & v
\end{array}
\]
defines a right inverse to the connecting homomorphism $\delta\colon \R(Z)_{(-a,b)} \to N_{(-a,b)}$.
\end{proposition}
Note that if $-a<0$ then necessarily $u^{(2)}=0$, and the lifting map $\epsilon$ is unique. If however $u^{(2)}\neq 0$, then 
it is not; we make a ``symmetric'' choice for our map $\epsilon$. 
\begin{remark}
    We note that the choice of $\epsilon$ is compatible with the $R$-module structures on $\bigoplus_{b\in \Z} N_{(-a,b)}$ and $\bigoplus_{b \in \Z} \R(Z)_{(-a,b)}$ for each $a \in \Z$. We can thus view $\epsilon$ as an $R$-module map from $\bigoplus_{b\in \Z} N_{(-a,b)}$ to $\bigoplus_{b \in \Z} \R(Z)_{(-a,b)}$ for each $a$, which we will use to prove the next Proposition.
\end{remark}

\subsection{Finding some sections in the Cox ring}
In this subsection, we complete Step~(3) in our procedure for all degrees $(-a,b)$ with $-1 \leq -a \leq d-2$, under the only assumption that $Z$ is nonsingular. We find $d$ new sections in the Cox ring of~$Z$ in degree~$(-1,e)$, and use these to describe $\R(Z)_{(-a,b)}$ for all degrees $(-a,b)$ with $-a \geq -1$.  

\begin{proposition}\label{prop:z_k}
\begin{enumerate} \item[(i)] There exists sections $z_1, \dots, z_d \in \R(Z)_{(-1,e)}$, linearly independent over $R$, that satisfy the $(d+1)$ equations
\begin{equation}
    x_1 z_1 + g_0 = x_1 z_{2} + g_{1} - x_0 z_{1} = \dots = x_1 z_{d} + g_{d-1} - x_0 z_{d-1} = g_d - x_0 z_d = 0.
    \label{eq:coxz}
\end{equation}
\item[(ii)]The subalgebra $\R(Z)_{z_k}$ of $\R(Z)$ generated by $x_0, x_1, y_0, \dots, y_n, z_1, \dots, z_d$ contains $\R(Z)_{(-a,b)}$ for all $-a \geqslant -1, b \in \Z$.\end{enumerate}
\label{prop:coxz1}
\end{proposition}

\begin{proof}
By Proposition~\ref{prop:ker}(i) we have
\begin{equation}\label{kerdecomp}
\bigoplus_{b\in \Z} N_{(-1,b)} = \bigoplus_{k = 1}^d R[e]\cdot \frac{1}{x_0^{d-k+1} x_1^k} \cong R[e]^d.
\end{equation}
For $1 \leq k \leq d$, let $f_k = \frac{1}{x_0^{d-k+1}x_1^k}$ be the $k^{\text{th}}$ basis vector in the decomposition (\ref{kerdecomp}). The short exact sequence~\eqref{seq:short} implies $\delta$ and $\epsilon$ are inverse isomorphisms for $a = -1$. Thus setting $z_k = \epsilon f_k$ gives us $d$ $R$-linearly independent sections $z_1, \dots, z_d \in \R(Z)_{(-1,e)}$. Following the definition of $\epsilon$, each $z_k$ is given by the following $0$-cochain $(z_{k,ij})$ on $X$:

\begin{equation*}
z_{k,ij} =
 \left\{
  \begin{array}{ll}
			 \sum_{c-k \geqslant 0} g_c \frac{x_1^{c-k}}{x_0^{c-k+1}} & \mbox{if } i = 0, \\
			 -\sum_{c-k < 0} g_c \frac{x_0^{k-c-1}}{x_1^{k-c}} & \mbox{if } i = 1. \\
		\end{array}
		\right.
\end{equation*}
Since for any choice of $(i,j)$ we have $x_1 z_{k+1,ij} + g_k - x_0 z_{k,ij} = 0$, we conclude that, for $k = 0,1,\dots,d$,
\begin{equation*}
    x_1 z_{k+1} + g_k - x_0 z_{k} = 0,
\end{equation*}
where the undefined variables $z_0$ and $z_{d+1}$ are set to be zero. This proves Proposition \ref{prop:coxz1}(i). 


For part (ii), note that $\iota^*\R(X)$ is generated as an algebra by the sections $x_0,x_1,y_0,\dots,y_n$ since this is true for $\R(X)$. To prove the statement it thus suffices to prove that each section in $\epsilon(N_{(-a,b)})$, for each $-a \geq -1$, is in $\R(Z)_{z_k}$. Part (i) covers the $a=-1$ case. Proposition~\ref{prop:ker} gives that $N_{(-a,b)}= 0$ for $-a > d-2$ and that
\begin{equation*}
\bigoplus_{b\in \Z} N_{(-a,b)} = \bigoplus_{k = 1}^{a+d-1} R[e]\cdot \frac{1}{x_0^{a+d-k} x_1^k} \cong R[e]^{a+d-1}.
\end{equation*}
for $0 \leq -a \leq d-2$. Fix $a$ in this latter range and for each $1 \leq k \leq a+d-1$ set $f_{-a,k} = \frac{1}{x_0^{a+d-k}x_1^k}$. These $f_{-a,k}$ give an $R$-basis for $\bigoplus_{b \in \Z} N_{(-a,b)}$. Since $\delta$ and $\epsilon$ are $R$-module maps, we only need to prove that each $v_{-a,k} = \epsilon f_{-a,k}$ is in $\R(Z)_{z_k}$ to conclude our proof. Computing the $0$-cochains associated to the $v_{-a,k}$ and comparing to the $0$-cochains associated to the generators of $\R(Z)_{z_k}$, we find for example the equations
\[
v_{-a,k} = z_k x_0^{1-a} - \frac{1}{2} \sum_{0 \leq c-k \leq -a} g_c x_0^{-a-c+k} x_1^{c-k}
\]
for each $0 \leq -a \leq d-2$ and $1 \leq k \leq a+d-1$. Then each $v_{-a,k}$ is in $\R(Z)_{z_k}$ and we are done.
\end{proof}

We deduce

\begin{corollary} \label{cor:partialdescription} Let
\[T = \C[X_0,X_1,Y_0,\dots,Y_n,Z_1,\dots,Z_d]\]
be the $\Z^2$-graded $\C$-algebra with the generators having degrees $(1,0)$, $(0,1)$ and $(-1,e)$, respectively. Let $I\lhd T$ be the ideal 
\[ I = \langle X_1 Z_1 + g_0(Y_i), X_1 Z_2 + g_1(Y_i) - X_0 Z_1,\ldots, g_d(Y_i) - X_0 Z_d\rangle\]
with generators corresponding to the $d+1$ equations in (\ref{eq:coxz}).
Then the algebra homomorphism 
\[ \phi : T \rightarrow \R(Z)
\]
defined by $X_i\mapsto x_i$, $Y_i\mapsto y_i$, $Z_i\mapsto z_i$ has kernel $K = \ker \phi$ containing~$I$, inducing isomorphisms $(T/I)_{(a,b)}\cong \R(Z)_{(a,b)}$ whenever $a \geqslant -1$.
\end{corollary}
\begin{proof}
By Proposition \ref{prop:coxz1}, the ideal $I$ is contained in $K = \ker \phi$. To prove the corollary, it suffices to show that $(K/I)_{a,b} = 0$ for $a \geqslant -1$.
Firstly, let $r' = g_0(Y_j)X_0^d + \dots + g_d(Y_j)X_1^d$, which maps via $\phi$ onto the defining equation $r$ of $Z$, so that $r' \in K$. Since
\begin{align*}
    r = &X_0^d \left(X_1 Z_1 + g_0(Y_j)\right) + X_0^{d-1} X_1 \left(X_1 Z_2 + g_1(Y_j) - X_0 Z_1\right) + \\ &\dots + X_0 X_1^{d-1} \left(X_1 Z_{d-1} + g_{d-1}(Y_j) - X_0 Z_d\right) + X_1^d \left(g_d(Y_j) - X_0 Z_d\right),
\end{align*}
we have $r' \in I$. Let $f = f(X_i,Y_j,Z_k)\in K$ be a homogeneous degree $(a,b)$ polynomial, for $a \geqslant -1$. We prove that the image $\overline{f}$ of $f$ in $K/I$ is zero by considering cases in $a$. Note that any term in $f$ is a monomial in the $Y_j$ multiplied by a monomial $M = X_0^{\alpha_0} X_1^{\alpha_1} Z_1^{\beta_1}\dots Z_d^{\beta_d}$ such that $\alpha_0 + \alpha_1 - \beta_1 - \cdots - \beta_d = a$.

If $a > d-2$, the generating equations of $I$ allow us to rewrite such an $M$ purely as a polynomial in the $X_i$ and $Y_j$ modulo $I$. Without changing the class $\overline{f}$, we replace $f$ with a polynomial $f(X_i,Y_j)$ independent of the $Z_k$. Since $\R(Z)_{(a,b)} = \C[x_i,y_j]_{(a,b)} / r\cdot \C[x_i,y_j]_{(a-d,b-e)}$ by Proposition \ref{prop:ker}, that $f \in K$ implies that $r$ divides $f(x_i,y_j)$ in $\C[x_i,y_j]$, so that $r'$ divides $f$ and thus $f \in I$ and $\overline{f} = 0$.

If $-1 \leq a \leqslant d-2$, the procedure is similar. The generating equations of $I$ allow us to rewrite $M$ as a $\C[Y_j]$-linear combination of $X_0^{a+1}Z_1,\dots,X_0^{a+1}Z_{d-a-1}$ plus a polynomial independent of the $Z_k$ modulo $I$. We can thus replace $f$, without changing $\overline{f}$, with a polynomial of the form
\begin{equation*}
    f(X_i,Y_j,Z_k) = f_1(Y_j)X_0^{a+1} Z_1 + \dots + f_{d-a-1}(Y_j)X_0^{a+1}Z_{d-a-1} + p(X_i,Y_j).
\end{equation*}
Following the proof of Proposition \ref{prop:coxz1}, we have
\begin{equation*}
    \phi(f) = f_1(y_j) v_{a,1} + \dots + f_{d-a-1}(y_j) v_{a,d-a-1} + q(x_i,y_j)
\end{equation*}
for some $q \in \C[x_i,y_j]_{(a,b)}\cong\iota^*\R(X)_{(a,b)}$. But $q, v_{a,1}, \dots, v_{a,d-a-1}$ are $R$-linearly independent, and so $\phi(f) = 0$ implies $f_1= \dots = f_{d-a-1} = 0$, and in turn $q = p = 0$. Then $\bar f=0$ as required.
\end{proof}

Proposition \ref{prop:coxz1}(ii) says that any Cox ring generators other than $x_i, y_j, z_k$ must be in degree $(a,b)$ with $a \leq -2$. Corollary \ref{cor:partialdescription} tells us that any new relations between generators must also be in degree $(a,b)$ with $a \leq -2$.

\subsection{Multiplying sections}

Proposition~\ref{prop:coxz1} and Corollary~\ref{cor:partialdescription} give a full picture of $\R(Z)$ in degrees $(-a,b)$ with $-a \geq -1$, for an arbitrary set of defining polynomials $\{g_c\}$ as long as $Z$ is non-singular. The computation of the remaining sections was reduced in Proposition~\ref{prop:ker} to the description of the kernels of the maps $A_a$. For general $\{g_c\}$, we are unable to give a full description; even the structure of $\ker A_2$ is hard to describe explicitly in general. In this section, we prepare the ground for further computations by utilising the identifications of each $\ker A_a$ with a submodule of $R^{a+d-1}$ to interpret multiplication by $x_0, x_1$ as well as $z_1,\dots, z_d$ in a succinct way. 

First, consider the multiplication map
\[ \cdot x_i \colon \R(Z)_{(-a,b)} \to \R(Z)_{(-a+1,b)}. 
\]
Composing this map with our lifting map $\epsilon$ on the right, and the projection map $\delta$ on the left, we obtain a map
\[\tilde x_i : \ker A_{a} \rightarrow \ker A_{a-1}.\]
\begin{proposition}
    Suppose that $-a \leqslant -2$. For $i=0,1$, the map $\tilde x_i : \ker A_{a} \rightarrow \ker A_{a-1}$ truncates the $(a+d-1)$-tuple $(f_m)_{m=1}^{a+d-1}$ on the right, respectively on the left, by one term.
    \label{xmult}
\end{proposition}
\begin{proof}
Fix $-a \leqslant -2, b \in \Z$ and let $(f_m)_{m=1}^{a+d-1} \in \ker A_a$,
%Recall that $\delta w$ is given by the image of 
identified with $f = \sum_{m=1}^{a+d-1} \frac{f_m}{x_0^{a+d-m} x_1^m}$.
Using Section~\ref{sec:lifting}, we obtain the lifting $w=\epsilon f$ defined by the degree $(-a,b)$ $0$-cochain $(w_{ij})$ given by
\begin{eqnarray*}
w_{0j} & = & f_1 g_d \frac{x_1^{d-1}}{x_0^{a+d-1}} + (f_1 g_{d-1} + f_2 g_d)\frac{x_1^{d-2}}{x_0^{a+d-2}} + \dots + (f_1 g_1 + \dots + f_d g_d) \frac{1}{x_0^a}, \\
w_{1j} & =  &  -f_{a+d-1} g_0 \frac{x_0^{d-1}}{x_1^{a+d-1}} - (f_{a+d-2} g_{0} + f_{a+d-1} g_1)\frac{x_0^{d-2}}{x_1^{a+d-2}} - \dots - (f_a g_0 + \dots + f_{a+d-1} g_{d-1}) \frac{1}{x_1^a}.
\end{eqnarray*}
% If $-a = -2$, then 
If $^-f$ is identified with the element of $\ker A_{a-1}$ obtained from truncating $(f_m)_{m=1}^{a+d-1}$ on the right, namely then $(f_m)_{m=1}^{a+d-2}$, then repeating this procedure for $^-w = \epsilon (^-f)$, we notice that $x_0 w_{0j} = ^-w_{0j}$ for each $j$, thus $x_0 w$ and $^-w$ agree on the dense open sets $Z \cap W_{0j}$ and therefore $x_0 w = ^-w$. If $f^-$ is the truncation on the left with corresponding section $w^- = \epsilon (f^-)$, then similarly $x_1 w_{1j}$ and $w^-_{1j}$ agree on the sets $Z \cap W_{1j}$ thus also $x_1 w = w^-$. 
\end{proof}

We now turn our attention to multiplying by $z_1, \dots, z_d$. 

\begin{proposition}
    Suppose that $-a \leqslant -1$. For each $k = 1,\dots,d$, the map $\tilde z_k : \ker A_{a} \rightarrow \ker A_{a+1}$, corresponding to multiplication by $z_k$, is given as follows. Suppose that $(f_m)_{m=1}^{a+d-1} \in \ker A_a$. Then $\tilde z_k (f_m)_{m=1}^{a+d-1}$ is given by $(f'_m)_{m=1}^{a+d}$, where
    \begin{align*}
        f'_m =& f_m g_k + \dots + f_{m+d-k} g_d, \quad m = 1,\dots,a+k-1, \\
        f'_m =& -f_{m-k} g_0 - \dots - f_{m-1} g_{k-1}, \quad m = k+1,\dots, a+d.
    \end{align*}
    If $-a = -1$, then this is specifies all $f'_m$. If $-a < -1$ then, for $m = k+1, \dots, a+k-1$, the two given definitions of $f'_m$ agree, as $(f_m)_{m=1}^{a+d-1} \in \ker A_a$.
    \label{zmult}
\end{proposition}

\begin{proof}
    Fix $-a \leqslant -1$. One checks that the maps as defined are indeed $R$-module homomorphisms with image in $\ker A_{a+1}$. Suppose that $(f_m)_{m=1}^{a+d-1} \in \ker A_a$ and let $\tilde z_d (f_m)_{m=1}^{a+d-1} = (f'_m)_{m=1}^{a+d}$. Then using Proposition~\ref{xmult} for the map $\tilde x_0$, we obtain
    \begin{equation*}
        \tilde{x}_0 \tilde{z}_d (f_m)_{m=1}^{a+d-1} = (f'_m)_{m=1}^{a+d-1}.
    \end{equation*}
    But the first equation in (\ref{eq:coxz}) tells us that the map $\tilde x_0 \tilde z_d$ is the same as multiplication by $g_d$, giving $f'_m$ as defined in the Proposition for the range $m = 1, \dots, a+d-1$. Since $(f'_m)_{m=1}^{a+d} \in \ker A_{a+1}$, the element $f'_{a+d}$ is uniquely determined by $f'_{a},\dots,f'_{a+d-1}$ and must also be given as in the Proposition. With the result proved for $k=d$, we proceed with a downwards induction in $l$. Suppose the formula is true for $k = d-k'$ and let $\tilde z_{d-k'-1} (f_m)_{m=1}^{a+d-1} = (f'_m)_{m=1}^{a+d}$. By the corresponding equation in \eqref{eq:coxz} we have
    \begin{equation*}
        (f'_m)_{m=1}^{a+d-1} = \tilde{x}_0 \tilde{z}_{d-k'-1} (f_m)_{m=1}^{a+d-1} = (f_m g_{d-k'-1})_{m=1}^{a+d-1} +\tilde{x}_1 \tilde{z}_{d-k'}(f_m)_{m=1}^{a+d-1}.
    \end{equation*}
    Using Proposition~\ref{xmult} for $\tilde x_1$ and the induction hypothesis we obtain the claimed formula for $f'_m$ for $m = 1,\dots,a+d-1$. The corresponding formula for $f'_{a+d}$ again follows as this is uniquely determined by $f'_{a},\dots,f'_{a+d-1}$.
\end{proof}

\subsection{The main theorem}

In this subsection, we prove Theorem~\ref{thm:mainintro}. 

%In this subsection, we assume that $\{g_0, \dots, g_d\}$ form a regular sequence in $R$. Note that we must then have $1 \leqslant d \leqslant n$, since $R$ has projective dimension $n$. In this case, the generators known already generate the full Cox ring.

\begin{theorem}\label{thm:general} Fix $n \geq 3$ and $d,e \geq 1$. Let $Z \subset X$ is non-singular, defined by a degree $(d,e)$ equation $r \in \R(X)$. Suppose that $Z$ is regular with index sequence $\{i_0,\dots,i_l\}$, where $1 \leq l \leq d$. Then there are sections $w_k \in \R(Z)_{(i_{k-1} - i_k,e)}$ for each $1 \leq k \leq l$ satisfying the equations
    \begin{equation}\label{eq:coxgen}
    x_1^{i_1 - i_0} w_1 + g_{i_0} = x_1^{i_2 - i_1} w_2 + g_{i_1} - x_0^{i_1 - i_0} w_1 = \dots = g_{i_l} - x_0^{i_l - i_{l-1}} w_l = 0.
   \end{equation}
    Furthermore, these sections, along with the sections $x_0,x_1,y_0, \dots,y_n$ restricted from $X$, generate the Cox ring $\R(Z)$. The ideal of relations among these generators is generated by the equations \eqref{eq:coxgen}. More precisely, let $U = \C[X_0, X_1, Y_0, \dots, Y_n, W_1,\dots,W_l]$ be the free $\Z^2$-graded polynomial ring with $\deg(X_i) = (1,0), \deg(Y_j) = (0,1)$ and $\deg(W_k) = (i_{k-1} - i_k,e)$ for the relevant $i,j,k$, and $J\lhd U$ the homogeneous ideal
    \begin{equation*}
       J = \langle X_1^{i_1 - i_0} W_1 + g_{i_0}(Y_j), X_1^{i_2 - i_1} W_2 + g_{i_1}(Y_j) - X_0^{i_1 - i_0} W_1,  \dots , g_{i_l}(Y_j) - X_0^{i_l - i_{l-1}} W_l\rangle. 
    \end{equation*}
    Then there is a surjective homomorphism of $\Z^2$-graded algebras
    $\psi : U \rightarrow \R(Z)$, inducing an isomorphism $\R(Z) \cong U/J$. In particular, the hypersurface~$Z$ is a Mori Dream Space.
\end{theorem}

\begin{proof}
    The sections $w_k$ are immediately found from the defining equation of $Z$. Indeed, for each $1 \leq k \leq l$,
    \begin{equation}\label{eq:wkdef}        
w_k=  \frac{g_{i_k} x_0^{d-i_k} + \dots + g_{i_l}x_1^{d-i_k}}{x_0^{d-i_{k-1}}} = -\frac{g_{i_0} x_0^{i_{k-1}} + \dots + g_{i_{k-1}}x_1^{i_{k-1}}}{x_1^{i_{k}}} \in \R(Z)_{(i_{k-1}-i_k,e)}
    \end{equation}
    is defined globally on $Z$, and these sections satisfy the equations (\ref{eq:coxgen}) in Theorem~\ref{thm:general}.
    In terms of our identification in Proposition~\ref{prop:ker}, the section $w_k$ for $1 \leq k \leq l$ is associated to the element $(f_m)_{m=1}^{i_k - i_{k-1} + d - 1} \in \ker A_{i_k - i_{k-1}}$ with $f_m = 1$ if $m = i_k$ and all other $f_m$ equal zero. We thus have a well-defined map $\psi\colon U\to \R(Z)$ mapping $X_i, Y_j, W_k$ respectively to the sections $x_i, y_z, w_k\in\R(Z)$, respectively. 

    The sections $z_1,\dots, z_d\in \R(Z)$ from Proposition~\ref{prop:coxz1}(i) are in the image of $\psi$; namely, for each $1 \leq k \leq l$ and each $1 \leq m \leq i_k - i_{k-1}$, it is easy to check that 
    \begin{equation} \label{eq:zinxw}
    z_{i_{k-1} + m} = \psi(Z_{i_{k-1}+m}) \ \mbox{ with } \ Z_{i_{k-1}+m} =X_0^{m-1}X_1^{i_k - i_{k-1} - m} W_k.
    \end{equation}
 So $\im\psi$ contains $\R(Z)_{(-a,b)}$ for all $-a \leq -1$ by Proposition~\ref{prop:coxz1}(ii). Furthermore, the equations \eqref{eq:coxz} become exactly the equations \eqref{eq:coxgen} and any degree $(a,b)$ monomial of $U$ with $a \geq -1$ can be rewritten as a monomial in  $X_0, X_1,Y_0,\dots,Y_n,$ $Z_1,\dots,Z_d$. Thus, by Corollary~\ref{cor:partialdescription}, we only need to consider $\R(Z)$ in degrees $(-a,b)$ with $-a \leq -2$. Surjectivity of $\psi$ in degrees $(-a,b)$ with $-a \leq -1$ is given by Theorem~\ref{regularsequencekernels}. It remains to show that the induced surjective map $\R = U/J \rightarrow \R(Z)$ is in fact an isomorphism.
 
Since $\psi$ is surjective, $Z$ is a Mori Dream Space, and thus $\dim \R(Z) = \dim Z + \rho (Z) = n+2$ by \cite[Proposition 2.9]{MDS}. To conclude the proof of the theorem, it therefore suffices to show that the algebra $\R = U/J$ is integral of dimension $n+2$. 
%\atsushi{Though I mention the normality in the previous file, it does not seem necessary to show that $\psi$ is an isomorphism. Hence it might be better to delete ``normal and '' in this sentence and the end of the proof of this theorem ``It follows that ...''.}


Localising $\R$ at $X_0$, we have the equations
\[
W_l = \frac{g_{i_l}(Y_j)}{X_0^{i_l - i_{l-1}}}, \dots, W_1 = \frac{X_1^{i_2 -i_1}W_2 + g_{i_1}(Y_j)}{X_0^{i_1 - i_0}}
\]
in $\R_{X_0}$, implying the surjectivity of 
\[
\C[X_0,X_1,Y_0,\dots,Y_n]_{X_0} \rightarrow \R_{X_0}.
\]
Substituting the above equations into $X_0^{i_1 - i_0} W_1 + g_{i_0}(Y_j) = 0$ in $\R_{X_0}$ we obtain $\frac{r(X_i,Y_j)}{X_0^d} = 0$ and thus
\begin{equation}\label{eq:Rlocalised0}
    \R_{X_0} = \left(\C[X_0,X_1,Y_0,\dots,Y_n]/r(X_i,Y_j)  \right)_{X_0}.
\end{equation}
Similarly one obtains
\begin{equation}\label{eq:Rlocalised1}
    \R_{X_1} = \left(\C[X_0,X_1,Y_0,\dots,Y_n]/r(X_i,Y_j)  \right)_{X_1}.
\end{equation}
On the other hand, we see
\begin{equation}\label{eq:Rirrelevant}
    \R / \langle X_0,X_1 \rangle = \C[Y_0,\dots,Y_n,W_1,\dots,W_l]/\langle g_{i_0}(Y_j), \dots, g_{i_l}(Y_j)\rangle.
\end{equation}
Let $T_1 = \Spec (\R / \langle X_0, X_1 \rangle)$ and $T_2 = \Spec (\R / \langle Y_0, \dots,Y_n \rangle)$. Since $r(X_i,Y_j)$ is irreducible, the rings $\R_{X_0}, \R_{X_1}$ are integral domains of dimension $n+2$ by \eqref{eq:Rlocalised0}, \eqref{eq:Rlocalised1} respectively. Additionally, since $\{g_{i_0},\dots,g_{i_l}\}$ form a regular sequence, we have $\dim T_1 = n$ by \eqref{eq:Rirrelevant}. One sees that $\dim (T_2 \cap \Spec \R_{X_0}) = 2$ via
\[
T_2 \cap \Spec \R_{X_0} = \Spec( \C[X_0,X_1,Y_0,\dots,Y_n]_{X_0}/\langle r(X_i,Y_j),Y_0,\dots,Y_j\rangle ) = \Spec \C[X_0,X_1]_{X_0}.
\]
In the same way, $\dim (T_2 \cap \Spec \R_{X_1}) = 2$. Since
\[
T_2 = (T_1 \cap T_2) \cup (\Spec \R_{X_0} \cap T_2) \cup (\Spec \R_{X_1} \cap T_2),
\]
we conclude that $\dim T_2 \leq \max \{\dim T_1,2 \} = n$. We may now consider the morphism
\[
\Spec \R \backslash (T_1 \cup T_2) \rightarrow Z,
\]
which is a $(\C^*)^2$-bundle. This implies that $\Spec \R \backslash (T_1 \cup T_2)$ is smooth and irreducible of dimension $n+2$. But $J \lhd U$ is generated by $l+1$ elements, and thus each irreducible component of $\Spec \R$ has dimension at least $n+2$. 
%Since $T_1 \cup T_2$ is in the closure of $\Spec \R \backslash (T_1 \cup T_2)$ and $\dim (T_1 \cup T_2) = n < n+2$, 
Thus $T_1 \cup T_2$ is not an irreducible component of $\Spec \R$, as $\dim T_1 \cup T_2 \leq n$.
Hence $T_1 \cup T_2$ is in the closure of $\Spec \R \backslash (T_1 \cup T_2)$, and 
we have that $\Spec \R$ is irreducible of dimension $n+2$. 
%It follows that $J$ is a complete intersection ideal and that $\Spec \R$ is smooth in codimension one, indeed it is smooth away from $T_1 \cup T_2$. Since $J$ is a complete intersection ideal and $\Spec \R$ is smooth in codimension one, it follows that $\R$ is integral \balazs{Do we need a reference here?}, and we are done.
It follows that $J$ is a complete intersection ideal and hence $\Spec \R$ is Cohen-Macaulay.
On the other hand, $\Spec \R$ is smooth in codimension one, indeed it is smooth away from $T_1 \cup T_2$.
Hence $\Spec \R$ is normal by Serre criterion.
In particular, $\Spec \R$ is integral, and we are done.
\end{proof}

%\begin{theorem}
    %Suppose that $\{g_0, \dots, g_d\}$ form a regular sequence in $R$. The sections $x_0, x_1, y_0, \dots, y_n$, restricted from $X$, along with the sections $z_1, \dots, z_d$ defined in Proposition~\ref{prop:coxz1}(i), generate the Cox ring $\R(Z)$. The ideal of relations between these generators is generated by the $(d+1)$ equations in (\ref{eq:coxz}). More precisely, let $T, I$ and $\phi$ be as in Corollary \ref{cor:partialdescription}. Then $\phi$ induces an isomorphism
    %$\bar\phi\colon T/I \to \R(Z)$. In particular, $Z$ is a Mori Dream Space.
    %\label{thm:regseq}
%\end{theorem}

%\begin{proof} We first show that the known sections generate $\R(Z)$, in other words the surjectivity of
%\[ \phi : \C[X_0,X_1,Y_0,\dots,Y_n,Z_1,\dots,Z_d] \rightarrow \R(Z). 
%\]
%From Corollary \ref{cor:partialdescription}, we know this surjectivity in degrees $(-a,b)$ with $-a \geq -1$. On the other hand, by Theorem~\ref{regularsequencekernels} below, for each $a \geq 2$, $\ker A_a$ is generated as an $R$-module by the elements corresponding to the degree $a$ homogeneous monomials in $z_1,\dots,z_d$. This proves the surjectivity of $\phi$ in degrees $(-a,b)$ with $-a\leq -2$ also. 

%To determine the full ideal of relations, we continue as in the proof of 
%Corollary~\ref{cor:partialdescription}. Let $K = \ker \phi$; clearly $I\subset K$. Suppose that $f = f(X_i, Y_j,Z_k) \in K$ has degree $(-a,b)$ with $-a \leq -2$. As in Corollary~\ref{cor:partialdescription}, any term in $f$ is a monomial in the $Y_j$ multiplied by a monomial $M = X_0^{\alpha_0} X_1^{\alpha_1} Z_1^{\beta_1}\dots Z_d^{\beta_d}$, with $\alpha_0 + \alpha_1 - \beta_1 - \cdots - \beta_d = -a$. The generating equations of $I$ allow us to rewrite such an $M$ purely as a polynomial in the $Y_j$ and $Z_k$ modulo $I$. Without changing the class $\overline{f} \in K/I$, we replace $f$ with a polynomial $f(Y_j,Z_k)$ independent of the $X_i$. 

%From $\phi(f)=0$, we see that $f(y_j,z_k)$ lies in the module of relations between the $y_j$ and $z_k$, which is fully described in Theorem~\ref{regularsequencekernels}. Thus $f(Y_j,Z_k)$ is an $R$-linear combination of the equations
%\begin{equation*}
    %g_k(Y_j) (Z_{l+1} Z_n - Z_l Z_{n+1}) - g_l(Y_j) (Z_{k+1} Z_n - Z_k Z_{n+1}) + g_n(Y_j) (Z_{k+1} Z_l - Z_k Z_{l+1})
%\end{equation*}
%for $1 \leq k < l < n \leq d$. For each $1 \leq k < l < n \leq d$, we have
%\begin{align*}
    %&g_k(Y_j) (Z_{l+1} Z_n - Z_l Z_{n+1}) - g_l(Y_j) (Z_{k+1} Z_n - Z_k Z_{n+1}) + g_n (Z_{k+1} Z_l - Z_k Z_{l+1}) \\
    %= &(Z_{l+1} Z_{n} - Z_{l} Z_{n+1})(X_1 Z_{k+1} + g_k(Y_j) - X_0 Z_{k}) \\
    %- & (Z_{k+1} Z_n - Z_k Z_{n+1})(X_1 Z_{l+1} + g_l(Y_j) - X_0 Z_{l}) \\
    %+ &(Z_{k+1} Z_{l} - Z_{k} Z_{l+1})(X_1 Z_{n+1} + g_n(Y_j) - X_0 Z_{n}).
%\end{align*}
%Thus $f \in I$ and we are done.
%\end{proof}

\begin{proposition}\label{prop:converse}
    Fix $n \geq 3$ and $d,e \geq 1$. Let $Z \subset X$ is non-singular, defined by a degree $(d,e)$ equation $r \in \R(X)$, admitting index sequence $\{i_0,\dots,i_l\}$. Let $U = \C[X_0, X_1, Y_0, \dots, Y_n, W_1,\dots,W_l]$ be the free $\Z^2$-graded polynomial ring with $\deg(X_i) = (1,0), \deg(Y_j) = (0,1)$ and $\deg(W_k) = (i_{k-1} - i_k,e)$ for the relevant $i,j,k$, and $J\lhd U$ the homogeneous ideal
    \begin{equation*}
       J = \langle X_1^{i_1 - i_0} W_1 + g_{i_0}(Y_j), X_1^{i_2 - i_1} W_2 + g_{i_1}(Y_j) - X_0^{i_1 - i_0} W_1,  \dots , g_{i_l}(Y_j) - X_0^{i_l - i_{l-1}} W_l\rangle.
    \end{equation*}
    If there is an isomorphism $\R(Z) \cong \R$ of $\Z^2$-graded algebras, where $\R = U/J$, then the sequence $g_{i_0},\dots,g_{i_l}$ is regular in $\C[y_0,\dots,y_n]$. In other words, $Z$ is regular with index sequence $\{i_0,\dots,i_l\}$.
\end{proposition}

\begin{proof}
    By degree considerations, the isomorphism $\R \cong \R(Z)$ sends the $\C$-span of $X_0, X_1$ to the $\C$-span of $x_0,x_1$. By \cite[Proposition 2.7]{MDS}, we conclude $X_0, X_1$ is a regular sequence in $\R$, and thus
    \[
    \dim \R / \langle X_0, X_1 \rangle = \dim \R -2 = \dim \R(Z) -2 = n,
    \]
    where again the final equality uses $Z$ being a Mori Dream Space. Using \eqref{eq:Rirrelevant} as found in the proof of the previous Theorem, we conclude that $\dim \C[Y_0,\dots,Y_n]/\langle g_{i_0}(Y_j), \dots, g_{i_l}(Y_j)\rangle = n-l$, implying that $g_{i_0},\dots,g_{i_l}$ is regular in $\C[y_0,\dots,y_n]$, as claimed.
\end{proof}

\section{Examples}

We look at some examples of our results of geometric interest. 

\begin{example} \label{ex:def1} Let $n=3$, $d=2$ and choose three general polynomials $g_0,g_1,g_2\in \C[y_0, y_1, y_2,y_3]$ of degree $e=4$, forming a regular sequence. Consider the family of varieties $q\colon {\mathcal Z} \to \A^1_t$ defined by
\[{\mathcal Z} = \left\{g_0 x_0^2 + t g_1 x_0x_1 + g_2 x_1^2 =0 \right\}\subset \p^1\times \p^3\times \A^1_t.
\]
For every $t\in\A^1$, the hypersurface fibre $Z_t\subset \p^1\times \p^3$ of the family $q\colon {\mathcal X} \to \A^1_t$ is a smooth Calabi--Yau threefold. Let $Z_t\stackrel{f_t}\longrightarrow {\bar Z}_t \stackrel{g_t}\longrightarrow \p^3$ be the Stein factorization of the second projection $p_2$. The map $g_t\colon {\bar Z}_t \to \p^3$ is a double cover in all cases, ramified over the divisor
\[ D_t = \{t^2g_1^2-4 g_0g_2=0\}\subset\p^3
\]
that is singular along the degeneracy locus 
\[Y_t=\{g_0=tg_1=g_2=0\}\subset\p^3.\]
For $t\neq 0$, $Y_t\subset\p^3$ is a set of $64$ points, and the map $f_t\colon Z_t\to {\bar Z}_t$ is a small contraction, contracting a set of $64$ rational curves to nodes. For $t=0$ on the other hand, the map $Z_0\to {\bar Z}_0$ is a divisorial contraction, contracting a divisor $E\subset Z_0$ to a genus-$33$ curve $C\subset {\bar Z}_0$ isomorphic 
to~$Y_0\subset \p^3$. As already observed by~\cite{constantinetal}, see in particular~\cite[3.3.2-3.3.3]{constantin}, for numbers of global sections, the Cox ring detects this change in birational behaviour. 
\begin{theorem}\label{thm_cy} For $t\neq 0$, the $\Z^2$-graded Cox ring $\R(Z_t)$ of the Calabi--Yau hypersurface $Z_t\subset \p^1\times \p^3$ can be presented as
\[ \R(Z_t)\cong k[x_0,x_1,y_0,y_1,y_2,y_3,z_1,z_2]/\langle x_1 z_1 + g_0, x_1 z_2 + g_1 - x_0 z_1, g_2 - x_0 z_2\rangle,\]
with variables of bidegrees $(1,0), (0,1)$ and $(-1,4)$ respectively. 
For $t=0$, we have 
\[\R(Z_0)\cong \C[x_0,x_1,y_0,y_1,y_2,y_3,w]/\langle x_1^2 w + g_0, g_2 - x_0^2 w\rangle,\]
with variables of bidegrees $(1,0), (0,1)$ and $(-2,4)$ respectively. 
In particular, every member of the family is a Mori dream space, with a complete intersection Cox ring, but the effective cone and Cox ring jump discontinuously in the family.  
\end{theorem}
\begin{proof} By the assumptions, for $t\neq 0$, the hypesurface $Z_t$ is regular with full index sequence $\{0,1,2\}$. On the other hand, the hypersurface $Z_0$ is regular, admitting index sequence $\{0,2\}$. Hence the statements follow from Theorem~\ref{thm:mainintro}. 
\end{proof}
Antonio Laface has informed us that the Cox rings in these examples can also be computed using the method of~\cite{lafaceetal}. 
\end{example}

We introduce one further, singular, member of this deformation family of varieties with interesting behaviour.

\begin{example} \label{ex:det} Choose general linear, respectively cubic polynomials $a_0,a_1,a_2\in R_1$ and 
$b_0,b_1,b_2\in R_3$. Consider the determinantal hypersurface
\[
Z = \left\{\begin{vmatrix}a_0 & a_1& a_2 \\ b_0 & b_1& b_2 \\ x_0^2 & x_0x_1 & x_1^2\end{vmatrix}=0\right\}
=\left\{g_0x_0^2  +g_1x_0x_1 +g_2 x_1^2 = 0\right\} \subset\p^1\times \p^3,\]
with $g_0=a_1b_2-a_2b_1, g_1=a_2b_0-a_0b_2, g_2=a_0b_1-a_1b_0$. 
Then the degeneracy locus is the curve
\[
Y = \left\{{\mathrm{rk}}\begin{pmatrix}a_0 & a_1& a_2 \\ b_0 & b_1& b_2 \end{pmatrix}\leq 1\right\}
 \subset \p^3,\]
a smooth space curve of genus $21$ and degree $13$. Its ideal $I_Y\lhd R$ has a resolution~\eqref{eq:sequence_for_T} in Hilbert--Burch form
 \begin{equation}0 \longrightarrow R^{2}\stackrel{B}\longrightarrow R^{3}\stackrel{A}\longrightarrow R \longrightarrow R/I_Y\longrightarrow 0.\end{equation}
Here, $B=\begin{pmatrix}a_0 & a_1& a_2 \\ b_0 & b_1& b_2 \end{pmatrix}^t$,  
and $A=\bigwedge^2 B = (g_0, g_1, g_2)^t$. As a determinantal hypersurface, $Z$ is singular along the locus given by the $2\times 2$ minors of its defining matrix, which gives the locus 
\[
\mathop\mathrm{Sing} Z=\left\{g_0=g_1=g_2=a_1^2-a_0a_2=x_0a_1-x_1a_0=x_0a_2-x_1a_1=0\right\}\subset \p^1\times \p^3,
\]
a set of~$26$ isolated ordinary double points all lying on the ruled surface $E\subset Z$. Blowing up this ruled surface, a Weil divisor through each of the ODP's, gives a small resolution $\widetilde Z\rightarrow Z$, a smooth Calabi--Yau model. 

Since $Z\subset X=\p^1\times\p^3$ has isolated nodal singularities, the restriction map $\pic(X)\to\pic(Z)$ is still an isomorphism. 
%if pushed for a proof, https://www.zib.de/userpage//lindner/leiden.pdf and	https://arxiv.org/pdf/1610.04077 would be useful
However, $Z$ is not $\Q$-Cartier, so the ring $\R(Z)$ as defined above only contains sections of Cartier divisors.

By Proposition~\ref{prop:ker}(ii), the columns of $B$ give us elements $f_1\in N_{(-2,5)}$ and $f_2\in N_{(-2,7)}$ that together generate the $R$-module $\bigoplus_{b \in \Z} N_{(-2,b)}$. Let $u=\epsilon f_1 \in \R(Z)_{(-2,5)}$ and $w =\epsilon f_2\in \R(Z)_{(-2,7)}$ be the lifts of these module generators to elements of the Cox ring as in Proposition~\ref{prop:epsilon}. By Proposition~\ref{xmult}, these satisfy equations
\begin{align*}
    x_0 u = a_1 z_1 + a_2 z_2, \quad x_1 u = a_2 z_1 + a_3 z_2, \\
    x_0 w = b_1 z_1 + b_2 z_2, \quad x_1 w = b_2 z_1 + b_3 z_2.
\end{align*}
Furthermore, by comparing $\delta(u), \delta(w),\delta(z_1^2),\delta(z_1 z_2),\delta(z_2^2)$ we have
\begin{equation*}
    z_1^2 + b_2 u - a_2 w = z_1 z_2 - b_1 u + a_1 w = z_2^2 + b_0 u - a_0 w = 0.
\end{equation*}
Consider the free $\Z^2$-graded $k$-algebra \[
    T = \C[X_0,X_1,Y_0, \dots,Y_n,Z_1,Z_2,U,W]
    \]
    with generators in degrees $(1,0), (0,1), (-1,4), (-2,5)$ and $(-2,7)$ respectively. Let $I \lhd T$ be the ideal 
    \begin{gather*}
       I= \langle X_1 Z_1 + g_0(Y_i), \quad X_1 Z_2 + g_1(Y_i) - X_0 Z_1, \quad g_2(Y_i) - X_0 Z_2,\\
        X_0 U - a_1(Y_i) Z_1 - a_2(Y_i) Z_2, \quad X_1 U - a_2(Y_i) Z_1 - a_3(Y_i) Z_2, \\
    X_0 W - b_1(Y_i) Z_1 + b_2(Y_i) Z_2, \quad X_1 W - b_2(Y_i) Z_1 - b_3(Y_i) Z_2,\\
    Z_1^2 + b_2(Y_i) U - a_2(Y_i) W, \quad Z_1 Z_2 - b_1(Y_i) U + a_1(Y_i) W, \quad Z_2^2 + b_0(Y_i) U - a_0(Y_i) W\rangle.
    \end{gather*}
    Then our observations so far prove the existence of an algebra homomorphism 
    \[
    \phi: T/I \rightarrow \R(Z),
    \]
 to the (Cartier) Cox ring  $\R(Z)$ of $Z$ that induces isomorphisms $(T/I)_{(a,b)} \cong \R(Z)_{(a,b)}$ whenever $a \geq -2$.


%\begin{conjecture} The map $\phi$ of Proposition~\ref{prop:linsyzpartialdescrip} is surjective with kernel $I$, giving an isomorphism $\R(Z)\cong T/I$. In particular, $\R(Z)$ is finitely generated.\end{conjecture} 
We sketch an argument provided to us by Antonio Laface that proves that for a general determinantal hypersurface~$Z$, the map $\phi$ 
% Proposition~\ref{prop:linsyzpartialdescrip} 
is an isomorphism. By~\cite[Thm.1]{lafaceetal}, the Cox ring $\R(Z)$ is the intersection of certain localizations of quotients
of~$\R(\p^1\times\p^3)$. This intersection can be computed using the ideas of~\cite[Cor.2.4]{lafaceetal}, which gives the result that~$\phi$ is surjective under certain dimension and saturation conditions. The latter can be checked for general $a_i,b_j$ using computer algebra.
\end{example} 

\begin{example}\label{ex:linear}
    Consider the hypersurface 
    \[Z=\{
    y_0 x_0^4 + (-\lambda_0 y_0 -\lambda_1 y_1 -\lambda_2 y_2 -\lambda_3 y_3)x_0^3 x_1 + y_1 x_0^2 x_1^2 + y_2 x_0 x_1^3 + y_3 x_1^4 = 0\}\subset \p^1 \times \p^3
    \]
    of bidegree $(d,e) = (4,1)$ with non-zero scalars $\lambda_j\in\C$, noting that a general degree $(4,1)$ hypersurface in $\p^1 \times \p^3$ has this form after a coordinate transformation.
    As the five linear forms are linearly dependent, by Theorem~\ref{thm:mainintro}(ii), the Cox ring is not of the form~\eqref{presentationR}, contrary to the statement \cite[Thm.1.1(iv)]{ottem}. We investigate the Cox ring of this particular family in this example. 

%\atsushi{I think that the following gives a proof of the conjecture.I have written this version a bit more carefully than the previous file, but please point out if there are any parts that are unclear.}

The kernel of $A_2$ as defined in Proposition~\ref{prop:ker} contains $(\lambda_0,1,\lambda_1,\lambda_2,\lambda_3)$, defining a section $w \in \R(Z)_{(-2,1)}$,
which satisfies the equations
\begin{equation}\label{eq:linear}
       x_0 w = \lambda_0 z_1 + z_2 +\lambda_1 z_3 +\lambda_2 z_4, \quad x_1 w = z_1 + \lambda_1 z_2+ \lambda_2 z_3 +\lambda_3 z_4  .
\end{equation}
    These equations
    %\[
    %x_0 w = \lambda_0 z_1 + z_2 +\lambda_1 z_3 +\lambda_2 z_4, \quad x_1 w = z_1 + \lambda_1 z_2+ \lambda_2 z_3 +\lambda_3 z_4 
    %\]
    imply that for general choices of $\lambda_j$, 
    the subalgebra $\R'$ of $\R(Z)$ generated by $x_0, x_1,w,z_3,z_4$ contains $z_1, z_2$.
    %%$\R(Z)_{z_k}$
    %$z_1, z_2$. 
    %we have %$z_k = p_k(x_0,x_1,w,z_3,z_4)$ for $k =1,2$, where
    %\begin{align*}
    %p_1 = \frac{1}{1-\lambda_0 \lambda_1} \left(w(x_1 - \lambda_1 x_0) + (\lambda^2_1 - \lambda_2)z_3 + (\lambda_1 \lambda_2 - \lambda_3)z_4\right), \\
    %p_2 = \frac{1}{1-\lambda_0 \lambda_1} \left(w(x_0 - \lambda_0 x_1) + (\lambda_0 \lambda_2 - \lambda_1)z_3 + (\lambda_0 \lambda_3 - \lambda_2)z_4\right).
    %\end{align*}
%\balazs{
Note also that the relations
\[
x_1 z_1 + y_0=0, \quad x_1 z_3 + y_1 - x_0 z_2=0, \quad x_1 z_4 + y_2 - x_0 z_3=0,\quad y_3 - x_0 z_4=0
%x_1 p_1 + y_0=0, \quad x_1 z_3 + y_1 - x_0 p_2=0, \quad x_1 z_4 + y_2 - x_0 z_3=0,\quad y_3 - x_0 z_4=0
\]
hold in $\R(Z)$.  %, which means that the $y_i$ are not needed as generators either.
%}
Hence the subalgebra $\R'$ contains $y_0, y_1,y_2,y_3$ as well.
In fact, we can show that $\R(Z)=\R'$, which is a polynomial ring as follows.

\begin{proposition}\label{prop:linear}
         For general constants $\lambda_j\in\C$, the $\Z^2$-graded Cox ring $\R(Z)$ of the hypersurface $Z \subset \p^1 \times \p^3$ is isomorphic to the polynomial ring $\C[x_0,x_1,w,z_3,z_4]$.    
\end{proposition}

\begin{proof}
Let $r \in H^0(\Ox_{\p^1 \times \p^3}(4,1)) =H^0(\Ox_{\p^1}(4)) \otimes H^0( \Ox_{\p^3}(1)) $ be the section which defines $Z$.
Then $r$ induces an exact sequence
\begin{align*}
0 \to \Ox_{\p^1}(-4) \xrightarrow{r} H^0( \Ox_{\p^3}(1)) \otimes \Ox_{\p^1} \to \mathcal{E} \to 0
\end{align*}
on $\p^1$
and we have $Z =\p_{\p^1}(\mathcal{E})\subset \p^1 \times \p^3$.
We note that $\mathcal{E}$ is locally free of rank $3$
since $r$ is general.
Let $\mathcal{E} \simeq \Ox_{\p^1} (a) \oplus \Ox_{\p^1} (b) \oplus \Ox_{\p^1} (c)$ with $a+b+c=4$.
Since $\mathcal{E}$ is a quotient of a trivial bundle, we have $a,b,c \geq 0$.
If $a =0$, then $r$ must be contained in $H^0(\Ox_{\p^1}(4)) \otimes V'$ for
$V' =\ker H^0(\Ox_{\p^3}(1)) \to H^0(\Ox_{\p^1}(a))=\C$,
which contradicts the generality of $r$.
Hence $a,b,c \geq 1$ and $\mathcal{E} \simeq  \Ox_{\p^1} (1) \oplus \Ox_{\p^1} (1) \oplus \Ox_{\p^1} (2)$. 


Since the embedding $Z =\p_{\p^1}(\mathcal{E})\subset \p^1 \times \p^3$ over $\p^1$ is induced by 
$H^0( \Ox_{\p^3}(1)) \otimes \Ox_{\p^1} \twoheadrightarrow \mathcal{E} $,
the tautological line bundle $\Ox_{\mathcal{E}}(1)$ on $ \p_{\p^1}(\mathcal{E})$ is identified with $\Ox_Z(0,1)$,
where $\Ox_Z(a,b) \coloneqq \Ox_{\p^1\times \p^3}(a,b)|_Z $.
Since $\mathcal{E}(-1) \simeq 
\Ox_{\p^1} \oplus \Ox_{\p^1} \oplus \Ox_{\p^1} (1)$, 
the morphism 
%\balazs{Balazs: I would write this as
%\[\mu : Z=\p_{\p^1}(\mathcal{E}(-1)) \to P = \p H^0(\Ox_Z(-1,1)) \cong\p^3\]
%}
\[\mu : Z=\p_{\p^1}(\mathcal{E}(-1)) \to P = \p H^0(\Ox_Z(-1,1)) \cong\p^3\]
induced by $|\Ox_Z(-1,1)|$ is the blow-up along a line $l$.
We note that $P$ is different from the second factor of $\p^1 \times \p^3$.


Let 
\begin{align*}
D_i =(x_i=0), \quad D'_k=(z_k =0), \quad E=(w=0) 
\end{align*}
be divisors on $Z$.
Since $w$ is a non-zero section of $H^0( \Ox_Z(-2,1)) = H^0(\mathcal{E}(-2)) \simeq \C$,
$E$ is the exceptional divisor of the blow-up $\mu$.

Since $z_1,z_2,z_3,z_4$ are linearly independent over $R$ by Proposition \ref{prop:z_k},
so are over $\C$.
Hence $z_1,z_2,z_3,z_4$ form a basis of $ H^0(\Ox_Z(-1,1))  =H^0(\mathcal{E}(-1)) \simeq \C^4$.
By (\ref{eq:linear}), $x_0w,x_1w,z_3,z_4$ also form a basis of $ H^0(\Ox_Z(-1,1)) =H^0(P, \Ox(1))$. We regard $P=\p^3$ as a toric variety by this basis.
Let $H_i=(x_iw=0) \subset P$ be the torus invariant planes for $i=0,1$.
Since $\mu^*H_i=D_i+E$, the line $l =\mu(E)$ is $H_0 \cap H_1$, which is torus invariant.
Hence $\mu$ is a toric morphism and $Z$ is a toric variety as well.
Then torus invariant prime divisors of $Z$ are $D_0,D_1,E, D'_3,D'_4$.
Since the Cox ring of a smooth projective toric variety is a polynomial ring whose variables correspond to torus invariants prime divisors,
we conclude that $\R(Z) \simeq \C[x_0,x_1,w,z_3,z_4]$.
\end{proof}



%    \begin{conjecture}
 %       For general constants $\lambda_j\in\C$, the $\Z^2$-graded Cox ring $\R(Z)$ of the hypersurface $Z \subset \p^1 \times \p^3$ can be presented as
%        \[      \R(Z) \cong \C[x_0,x_1,y_0,\dots,y_n,w,z_3,z_4]/J,       \]
 %       where $J$ is the ideal 
  %      \begin{align*}
   %         J = \left\langle x_1 p_1 + y_0, \quad x_1 p_2 + (-\lambda_0 y_0 -\lambda_1 y_1 -\lambda_2 y_2 -\lambda_3 y_3) - x_0 p_1, \right. \\ \left. x_1 z_3 + y_1 - x_0 p_2, \quad x_1 z_4 + y_2 - x_0 z_3,\quad y_3 - x_0 z_4\right\rangle.
    %    \end{align*}
   % \end{conjecture}
    %More generally, for a general hypersurface $Z \subset \p^1 \times \p^n$ of bidegree $(d,1)$ with $d > n$, the linear forms $g_0,\dots,g_d$ must be linearly dependent. Each linear dependency gives an element with scalar coordinates in $\ker A_2$, which gives a non-zero section in $\R(Z)_{(-2,1)}$, which cannot be generated by the sections $x_i, y_j, z_k$. %\R(Z)_{z_k}$
    %In particular, the Cox ring never take the form of Theorem~\ref{thm_main_intro}(i) in this range of parameters.
    
    %Even more generally for a hypersurface of degree $(d,e)$ in $\p^1 \times \p^n$ with $d,e \geq 1$, if the sequence $g_0,\dots,g_d$ is not regular in $R$, then there is at least one section on $Z$ in degree $(-2,b)$ for some $b \in \Z$ not contained in $\R(Z)_{z_k}$, and the Cox ring cannot take the form given in Theorem~\ref{thm_main_intro}.
\end{example}

As a last example, we discuss from our point of view a phenomenon in the birational geometry of the hypersurface $Z$ that formed a key part of the argument of~\cite{ottem}.

\begin{example}\label{ex:regcasemodel}
    Let $n \geq 3$ and $d,e \geq 1$. Choose general polynomials $g_0,\dots,g_d \in \C[y_0,\dots,y_n]$ of degree $e$, forming a regular sequence. Then
    \[
    \R(Z) \cong \C[x_0,x_1,y_0,\dots,y_n,z_1,\dots,z_d]/\langle x_1 z_1 + g_0, x_1 z_2 + g_1 - x_0 z_1,\ldots, g_d - x_0 z_d\rangle.
    \]
    Let $S_{Z^+}$ be the subalgebra of $\R(Z)$ generated by $y_0,\dots,y_n,z_1,\dots,z_d$. By Proposition~\ref{prop:generalWideal}, proved in the final section, we have an isomorphism
    \[
    S_{Z^+} \cong \C[y_0,\dots,y_n,z_1\dots,z_d]/J',
    \]
    where $J'$ is the ideal generated by the equations
    \begin{equation} \label{eq:secondmodel}
     g_{k} (z_{k'+1} z_{k''} - z_{k'} z_{k''+1}) - g_{k'} (z_{k+1} z_{k''} - z_k z_{k''+1}) + g_{k''} (z_{k+1} z_{k'} - z_k z_{k'+1}),
    \end{equation}
    for $0 \leq k < k' < k'' \leq d$. The ring $S_{Z^+}$, defined by the equations \eqref{eq:secondmodel}, is the homogeneous coordinate ring of a birational model $Z^+ \subset \p^d_{z_k} \times \p^n_{y_j}$ of $Z$; compare~\cite[Section 2]{ottem}.
\end{example}


\section{A problem in algebra}

%Let $R$ be an algebraically closed field of characteristic zero \balazs{Should I just make this $\C$ everywhere, I only use this to say that the Cox ring later has no zero divisors?} and $d \geq 0$.

Fix $d \geq 0$ and fix an admissible index sequence $\{i_0,\ldots, i_l\}$ for some $0 \leq l \leq d$; note that $d = 0$ is equivalent to $l = 0$ by our conventions. Let $S = \C[Y_{i_0}, \dots, Y_{i_l}]$ be the free commutative $\C$-algebra on $l+1$ generators. For each $a \geqslant 1$, define a map of free $S$-modules \[A_a : S^{a+d-1} \rightarrow S^{a-1}\] by the $(a-1) \times (a+d-1)$ matrix with $(s,t)$ entry equal to $Y_{i_k}$ if $t-s = i_k$ and $0$ otherwise. Denote $K_a=\ker A_a$ and let $K = \bigoplus_{a\geq 1} K_a$, which we regard as a graded $S$-module. Note that $A_1=0$ and so $K_1\cong S^d$, with the convention that $S^0 = 0$.

\begin{example}
    If $d = 5, l = 2$ and the degree sequence is $\{0,3,5\}$, then $A_3$ is given by
\begin{equation*}
    \begin{pmatrix}
   Y_{0} & 0 & 0 & Y_{3} & 0 & Y_{5} & 0 \\
   0 & Y_{0} & 0 & 0 & Y_{3} & 0 & Y_{5}
    \end{pmatrix}.
    \end{equation*}
\end{example}

We will mostly employ formal algebraic ideas in this section, but sometimes the use of the following auxiliary hypersurface will be helpful. Consider the non-singular hypersurface
    \[
    Z' = \{Y_{i_0} x_0^{d} + Y_{i_1} x_0^{d-i_1} x_1^{i_1} + \dots + Y_{i_l} x_1^d\} \subset \p_{x_i}^1 \times \p_{Y_{i_k}}^l.
    \]
We can repeat the arguments of Section 2 for this hypersurface. Its Cox ring contains sections $w_1,\dots,w_l\in\R(Z')$ as given in~\eqref{eq:wkdef}, and as in equations~\eqref{eq:zinxw}, contains the sections $z_1,\dots,z_d\in\R(Z')$ given in Proposition~\ref{prop:coxz1}. We also have an identification 
\begin{equation}\label{eq:KandZ'} K\cong \bigoplus_{-a \leq -1} \bigoplus_{b\in\Z} \R(Z')_{(-a,b)}\end{equation}
as graded $S$-modules via Proposition~\ref{prop:ker}.

\begin{definition}
    Fix $a \geq 1$.
\begin{itemize}
    \item[(i)] Define the operators $\tilde x_0, \tilde x_1\colon K_{a+1} \rightarrow K_a$ respectively as truncation by one entry on the right, respectively on the left.
    \item[(ii)] For each $1 \leq k \leq d$ we define $\tilde z_k: K_a \rightarrow K_{a+1}$ as follows. Suppose that $(f_m)_{m=1}^{a+d-1} \in \ker A_a$. Then $ \tilde z_k (f_m)_{m=1}^{a+d-1}$ is given by $(f'_m)_{m=1}^{a+d}$, where
    \begin{align*}
        f'_m =& f_m Y_k + \dots + f_{m+d-l} Y_d,& \quad m = 1,\dots,a+k-1, \\
        f'_m =& -f_{m-k} Y_0 - \dots - f_{m-1} Y_{k-1},& \quad m = k+1,\dots, a+d,
    \end{align*}
    where $Y$ variables of undefined index are set to be zero.
    \item[(iii)] For each $1 \leq k \leq l$, we define $\tilde w_k : K_{a} \rightarrow K_{a+i_k - i_{k-1}}$ by stipulating that
    \[
    \tilde x_0^{m-1} \tilde x_1^{i_k - i_{k-1}-m} \tilde w_k = \tilde z_{i_{k-1}+m}
    \]
    for each $1 \leq m \le i_k - i_{k-1}$.
\end{itemize}
    \label{operators1}
\end{definition}

\begin{proposition} \label{prop:opwelldefined}
The operators $\tilde x_i, \tilde z_k, \tilde w_k\colon K\to K$ in Definition~\ref{operators1} are well defined. 
\end{proposition}
\begin{proof} For $\tilde x_i$, the statement is clear. For $z_k$, when $a = 1$, then each $f'_m$ is specified by one formula. If $a > 1$, then for $m = k+1, \dots, a+k-1$, the two given definitions of $f'_m$ agree, since $(f_m)_{m=1}^{a+d-1} \in \ker A_a$. Similarly, it is easy to see directly that the operators $\tilde w_k$ are well defined. 

An alternative argument employs the auxiliary hypersurface $Z'$ defined above. 
The operators in Definition~\ref{operators1} correspond to multiplication by the corresponding sections $x_i, z_k, w_k\in \R(Z')$ by Propositions~\ref{xmult}--\ref{zmult} under the identification~\eqref{eq:KandZ'}, and are thus well defined.  
\end{proof}

Let $U = S[X_0, X_1, W_1,\dots,W_l]$ be the free graded $S$-algebra on $l+2$ generators, with $\deg(X_i)=1$ and $\deg(W_k) = i_{k-1}-i_k$ for relevant $i,k$. For a general graded $S$-module $U'$, we denote by $U'_{\leq -1}$ the graded submodule $\bigoplus_{a\leq -1} U'_{a}$. It is also useful to define the elements
\begin{equation} \label{eqn:ZinXWalg}
Z_{i_{k-1}+m} = X_0^{m-1}X_1^{i_k - i_{k-1}-m}W_k \in U_{-1},
\end{equation}
for $1\leq k \leq l$ and $1 \leq m \leq i_k - i_{k-1}$.

\begin{proposition} \label{prop:psialg}
    There is a map of $S$-modules
    \[
    \psi: U_{\leq -1} \rightarrow K
    \]
    defined as follows. Define, for each $1 \leq k \leq l$, the element $\psi(W_k) \in K_{i_k - i_{k-1}}$ to be $(f_m)_{m=1}^{i_k - i_{k-1} + d -1}$ with $f_m = 1$ if $m = i_k$ and all other $f_m$ equal 0. The grading on $U$ implies that any monomial $M$ in $U_{\leq -1}$ divides by $W_k$ for some $k$. If $M = X_0^{\alpha_0} X_1^{\alpha_1} W_1^{\beta_1}\dots W_l^{\beta_l}$ with say $\beta_k >0$, then we set
    \[
    \psi(M) = \tilde x_0^{\alpha_0} \tilde x_1^{\alpha_1} \tilde w_1^{\beta_1}\dots \tilde w_k^{\beta_k -1} \dots \tilde w_l^{\beta_l} \psi(W_k).
    \]
    Furthermore, the elements $\psi(Z_k) \in K_1$, for $1 \leq k \leq d$, give the standard basis of $K_1 \cong S^d$.
\end{proposition}

\begin{proof}
    To argue that $\psi$ is well defined, we need only argue the operators $\tilde x_i$ and $\tilde w_k$ all commute with each other, as well as establish that $\tilde w_k \psi(W_{k'}) = \tilde w_{k'} \psi(W_k)$ for each $1 \leq k, k' \leq l$. 

    It is possible to establish this claim directly algebraically, using the definitions of the operators. 
    This somewhat long-winded argument will feature in \cite{pollock}.
    %If $l = d$, one can show that the operators $\tilde z_k$ commute with each other directly from their algebraic definitions, as was done in an earlier version of this paper. The commutativity of the $\tilde z_k$ in the general setting follows immediately, for the formulas are the same with $Y_c$ set to be zero when $c \notin \{i_0,\dots,i_l\}$. The result follows with the $\tilde w_k$ operators as these are defined in terms of the $\tilde z_k$ operators and the truncation operators $\tilde x_0, \tilde x_1$, which clearly commute. This argument will feature in \cite{pollock}.

    To give a short argument instead, notice as above that the operators correspond to multiplication of sections in the Cox ring $\R(Z')$ of the auxiliary hypersurface $Z'$. Since multiplication of these sections is clearly commutative, all of the operators must commute with each other, and the extra property
    $\tilde w_k \psi(W_{k'}) = \tilde w_{k'} \psi(W_k)$ for each $1 \leq k, k' \leq l$ also holds.

    The final statement is easily checked.
\end{proof}

\begin{theorem}\label{thm:generalgebra}
    Let $J \lhd U$ be the homogeneous ideal defined by
    \begin{equation*}
       J = \langle X_1^{i_1 - i_0} W_1 +Y_{i_0}, X_1^{i_2 - i_1} W_2 + Y_{i_1} - X_0^{i_1 - i_0} W_1,  \dots , Y_{i_l} - X_0^{i_l - i_{l-1}} W_l\rangle. 
   \end{equation*}
    Then $\psi:U_{\leq -1} \rightarrow K$ is surjective with kernel $\ker \psi = J_{\leq -1}$. In particular there is an isomorphism of $S$-modules $K \cong (U/J)_{ \leq -1}$.
\end{theorem}

We prove Theorem~\ref{thm:generalgebra} using induction in $l$. If $l = 0$ then $d = 0$, the result is true since in this case $U_{ \leq -1} = K = 0$. We thus assume in the following that $l \geq 1$. For the induction step we need some Definitions and Lemmas.

\begin{definition}\label{def:inductionbar} Let $\bar S = \C[Y_{i_0},\dots,Y_{i_{l-1}}]$. For each $a \geq 1$, let 
\[\bar A_a : \bar S^{a+d'-1} \rightarrow \bar S^{a-1}\]
be the $\C$-module maps as defined at the start of this section, but for the case of $l$ variables with partition $0 = i_0 < \dots < i_{l-1} = d'$, with corresponding auxiliary hypersurface. Denote $\bar K_a=\ker \bar A_a$ and $\bar K = \bigoplus_{a \geq 1} \bar K_a$. Define the operators $\tilde \zeta_k\colon\bar K_a\to \bar K_{a+i_k - i_{k-1}}$ for $1 \leq k \leq l-1$ analogously to the $\tilde w_k$ in Definition~\ref{operators1}(iii). We abuse notation slighlty and still use $\tilde x_0$ and $\tilde x_1$ for the respective truncation operators. Define $\bar U = S[X_0,X_1,\bar \zeta_1,\dots,\bar \zeta_{l-1}]$, and define the ideal $\bar J \lhd \bar U$ and map $\bar \psi : \bar U_{\leq -1} \rightarrow \bar K$ analogously to $J$ and $\psi$ in Proposition~\ref{prop:psialg}.
\end{definition}

\begin{lemma}\label{lem:inductionw}
    Consider an element $\psi (W_{k_s} \cdots W_{k_1} W_{k_0}) \in K_t$, where no $k_i$ is equal $l$. Setting $Y_{i_{l}} = 0$ in the resulting element of $S^{t+d-1}$ and discarding the final $i_l - i_{l-1}$ components gives us an element of $\bar S^{t+d'-1}$. This coincides with $\psi(\bar \zeta_{k_s} \cdots \bar \zeta_{k_1} \bar \zeta_{k_0}) \in \bar K_{t}$.
\end{lemma}

\begin{proof}
    This follows immediately from the definitions.
\end{proof}

\begin{lemma}\label{lem:inddet}
    Suppose $u = (f_m)_{m = 1}^{a + d-1} \in K_a$. If $a > i_l - i_{l-1}$ then $u$ is determined by its first $a + i_{l-1} - 1$ entries. If $a \leq i_l - i_{l-1}$ then $u$ is determined by its first $a + i_{l-1} - 1$ entries up to free choices of the entries $f_{a+ i_{l-1}},\dots, f_{i_l}$.
\end{lemma}

\begin{proof}
    If $a > i_l - i_{l-1}$, then the first $i_l - i_{l-1}$ rows of $A_a$ give equations determining the remaining entries. If $a \leq i_l - i_{l-1}$, then the rows of $A_a$ determine each of the last $a-1$ entries in $u$ and the $a+i_{l-1}, \dots, i_l$ columns of $A_a$ are zero columns, giving the remaining free choices.
\end{proof}

\begin{proof}[Proof of Theorem~\ref{thm:generalgebra}]
    We know the result is true for $l = 0$, so suppose that $1 \leq l \leq d$. We first prove that $\psi$ is surjective. Suppose that $a \geq 1$ and $u = (f_m)_{m=1}^{a+d-1}\in K_a$ is given. Setting $Y_{i_l} = 0$ and discarding the last $i_l - i_{l-1}$ elements we obtain $\bar u = (\bar f_m)_{m=1}^{a+d'-1} \in \bar K_a$. By induction in $l$ we can find $p(X_0,X_1,\bar \zeta_1,\dots,\bar \zeta_{l-1}) \in \bar U_{\leq -1}$ such that $\bar \psi(p) = \bar u$. By Lemma~\ref{lem:inductionw} we have
    \[
    u' = u - \psi \left(p(X_0,X_1,W_1,\dots,W_{l-1})\right) = (f'_1 Y_{i_l},\dots,f'_{a+d'-1} Y_{i_l},F_{a+d'},\dots,F_{a+d-1}) \in K_a
    \]
    for some $f'_1,\dots,f'_{a+d'-1},F_{a+d'},\dots,F_{a+d-1} \in S$. We continue inductively in $a$.
    
    Suppose $1 \leq a \leq i_l - i_{l-1} = d-d'$. Let $v = (f'_1,\dots,f_{a+d'-1},0,\dots,0) \in K_1$, where the number of zeroes at the end is $d-d' -a +1$. By the final statement of Proposition~\ref{prop:psialg}, 
    \[v = \psi(f'_1 Z_1 + \dots + f'_{a+d'-1}Z_{a+d'-1})
    \]
    is in the image of $\psi$. Now consider $\tilde x_0^{d-d' -a +1} \tilde w_l (v) \in K_a$, which is equal
    \[
    (f'_1 Y_d,\dots,f'_{a+d'-1}Y_d,0,\dots,0,f''_{d+1},\dots,f''_{a+d-1}).
    \]
    Now $\tilde x_0^{d-d' -a +1} \tilde w_l (v)$, which is in the image of $\psi$ since $v$ is, and $u'$ have the same first $a+d'-1$ entries in common, thus the last $a-1$ entries of each must also agree by Lemma~\ref{lem:inddet}. Then
    \[
    u' - \tilde x_0^{d-d' -a +1} \tilde w_d (v) = (0,\dots,0,F_{a+d'},\dots,F_{d},0,\dots,0) \in K_a.
    \]
    But this is in turn equal to
    \[
    \psi\left((F_{a+d'} X_1^{d-d'-a} + \dots + F_{d} X_0^{d - d' -a})W_l\right),
    \]
    so that $u$ can be written in the image of $\psi$.

    If instead $a > d - d'$, then $(f'_1 Y_{i_l},\dots,f'_{a+d'-1} Y_{i_l}) \in K_{a-d+d'}$. Since $Y_{i_l}$ is not a zero divisor in $S$, we have $v = (f'_1,\dots,f'_{a+d'-1}) \in K_{a-d+d'}$. Since the first $a+d'-1$ coordinates of $u'$ and $\tilde w_l(v)$ are equal, we have $u' = \tilde w_l (v)$ by Lemma~\ref{lem:inddet}. By induction in $a$, $v$ lies in the image of $\psi$, and thus so does $u$.

    We now turn to the kernel of $\psi$. We can check that $J_{\leq -1} \subset \ker \psi$. Suppose that $p \in (\ker \psi)_{-a}$ for some $a \geq 1$. Consider $\bar p \in \bar U$ obtained from $p$ by setting $Y_{i_l} = W_l = 0$ and replacing each $W_k$ with $\bar \zeta_k$ for each $1 \leq k \leq l-1$. By Lemma~\ref{lem:inductionw}, $\bar p \in \ker \bar \psi$. By induction in $l$, we have $\bar p \in \bar J$, and thus $\bar p$ may be written as an $S$-linear combination of the equations
\[
X_1^{i_{k+1} - i_{k}} \bar \zeta_{k+1} + Y_{i_k} - X_0^{i_k - i_{k-1}} \bar \zeta_{k}, \quad 0 \leq k \leq l-1,
\]
where undefined indices are set to be zero. Write $\bar p$ as a $\bar U$-linear combination of these equations, except leaving the variable $\bar \zeta_l$ in where it is usually set it equal to zero. Let $p' \in U$ be the corresponding element of $J$ obtained from this linear combination by replacing $\bar \zeta_k$ with $W_k$ for each $1 \leq k \leq l$, so that $p' \in J_{\leq -1}$. We can replace $p$ with $p-p'$ without affecting whether it belongs to $J_{\leq -1}$. By construction, all terms of $p$ now divide by $Y_{i_l}$ or $W_l$. Adding a multiple of $Y_{i_l} - X_0^{i_l - i_{l-1}}W_l$, which is in $J$, we can replace $p$ with something dividing by $W_l$. Write $p = q W_l$.

If $a \leq i_l - i_{l-1}$ then $q \in U_{\geq 0}$. Consider the section \[q(x_0,x_1,w_1,\dots,w_l)\in\R(Z'),\] in the Cox ring of the auxiliary hypersurface $Z'$. The section $q(x_0,x_1,w_1,\dots,w_l)w_l \in \R(Z')$ is identified with $\psi(p) = 0$ in \eqref{eq:KandZ'} and is thus the zero section. Since $Z'$ is irreducible and normal, its Cox ring has no zero divisors, thus $q(x_0,x_1,w_1,\dots,w_l) = 0$ in $\R(Z')$. By Corollary~\ref{cor:partialdescription}, $q(X_0, X_1,W_1,\dots,W_l) \in \psi(J)$ . Then $p \in J_{\leq -1}$.

If $a > i_l -i_l$, then $0 = \psi (p) = \tilde w_l \psi(q)$. Now $\tilde w_l$ is an injective operator by Lemma~\ref{lem:inddet} and the fact it is multiplication by $Y_{i_l}$ on the first $a + d'-1$ terms. Thus $\psi(q) = 0$. By induction, $q \in J_{\leq -1}$ and thus $p\in J_{\leq -1}$. 
\end{proof}

We add a final Proposition, whose generalisation in Proposition~\ref{prop:generalWideal} is used for Example~\ref{ex:regcasemodel}.

\begin{proposition}\label{prop:universalWideal}
    Suppose that $l = d$ and $J, U$ are as in Theorem~\ref{thm:generalgebra}. Let $U' = S[W_1,\dots,W_d]$ be the subalgebra of $U$ generated by $W_1,\dots,W_d$. Let $J' \lhd U'$ be the homogeneous ideal generated by the equations
    \[
    Y_k (W_{k'+1} W_{k''} - W_{k'} W_{k''+1}) - Y_{k'} (W_{k+1} W_{k''} - W_k W_{k''+1}) + Y_{k''} (W_{k+1} W_{k'} - W_k W_{k'+1})
    \]
    for $0 \leq k < k' < k'' \leq d$. Then the natural map $\phi:U'_{\leq -1} \rightarrow (U/J)_{\leq -1} $ is surjective and has kernel $J'$.
\end{proposition}

\begin{proof}
    The equations correspond exactly to the images of the standard basis vectors in the degree $3$ map of the Koszul complex of the sequence $Y_0,Y_1,\dots,Y_d$ in $S$, and thus are contained in $J$. Alternatively, that $J' \subset J$ can be checked explicitly by writing the generators of $J'$ explicitly as $S$-combinations of the generators of $J$.
    
    We use induction in $d$. If $d = 0$ the result is immediate, so suppose that $d \geq 1$. Let $\bar U'$ and $\bar J'$ be the respective analogues of $U'$ and $J'$ with parameters $d-1 = l-1$. The equations in $J$ imply that any element in $(U/J)_{\leq -1}$ can be written as a polynomial in the $Y_j, W_k$, thus $\phi$ is surjective. Notice next that $(U/J)_{-1} \cong S^d \cong (U'/J')_{-1}$. Suppose that $p(Y_j,W_k) \in \ker \phi$ has degree $-a \leq -2$. Then $p(Y_j,W_k)$ is also in $\ker \psi = J$. As in the proof of Theorem~\ref{thm:generalgebra}, we can assume all terms of $p$ divide by $Y_d, W_d$. The equations of $J'$ can then be used to replace $p$ with a polynomial dividing by $W_d$, so that $p = qW_d$. Since $\tilde w_d$ is an injective operator, we conclude that $q \in \ker \phi$. Using induction in the degree~$a$, we have $q\in J'$ and thus $p \in J'$.
\end{proof}

\begin{theorem} \label{regularsequencekernels}
    Fix $d \geq 0$ and fix an admissible index sequence $\{i_0,\ldots, i_l\}$ for some $0 \leq l \leq d$. Let $S$ be a finitely generated graded $\C$-algebra with projective dimension $\kappa \geq l$. Suppose that $g_{i_0},\dots,g_{i_l}$ are homogeneous elements of degree $e$ forming a regular sequence in $S$. For each $a \geq 1$, define a map of free $S$-modules $A_a : S^{a+d-1} \rightarrow S^{a-1}$ by the $(a-1) \times (a+d-1)$ matrix with $(s,t)$ entry equal to $g_{i_k}$ if $t-s = i_k$ and $0$ otherwise. Let $K_a = \ker A_a$ and $K = \bigoplus_{a\geq 1} K_a$. Let $U = S[X_0,X_1,W_1,\dots,W_l]$ be the free $\Z^2$-graded $S$-algebra with $\deg_U(s) = (0,\deg_S(s)), \deg(X_i) = (1,0)$ and $\deg(W_k) = (i_{k-1}-i_k,e)$ for $s \in S$ and relevant $i,k$.     Then there is a surjective map
    \[
    \psi : U_{\leq -1} = \bigoplus_{-a \leq -1} \bigoplus_{b\in \Z} U_{(-a,b)} \rightarrow K,
    \]
whose kernel contains $J_{\leq -1}$, where $J \lhd U$ is the homogeneous ideal given by
    \[
       J = \langle X_1^{i_1 - i_0} W_1 + g_{i_0}(Y_j), X_1^{i_2 - i_1} W_2 + g_{i_1}(Y_j) - X_0^{i_1 - i_0} W_1,  \dots , g_{i_l}(Y_j) - X_0^{i_l - i_{l-1}} W_l\rangle.
    \]
\end{theorem}

\begin{proof}
    If $S = \C[Y_{i_0},\dots,Y_{i_1}]$ and $g_{i_k} = Y_{i_k}$ for each $k = 0,\dots, l$, then the statement is that of Theorem~\ref{thm:generalgebra} with refined degrees. All of the proofs given above can be modified to this more general setting. The operators in Definition~\ref{operators1} are defined in the same way with the $Y_{i_k}$ replaced by the $g_{i_k}$. Proposition~\ref{prop:psialg} holds as we symbolically only replace the $Y_{i_k}$ with $g_{i_k}$, preserving commutativity. In Definition~\ref{def:inductionbar}, we set $\bar S = S/(g_{i_l})$, which has projective dimension $\kappa -1 \geq l-1$. For Lemma \ref{lem:inductionw}, rather than setting $Y_{i_l} = 0$, we reduce modulo $g_{i_l}$. Lemma~\ref{lem:inddet} holds with regularity of the sequence $g_{i_0},\dots,g_{i_l}$ in $S$.

    Now the proof follows that of Theorem~\ref{thm:generalgebra}. We use induction in $(\kappa,l,a)$. If $\kappa = 0$ then $l = 0$, so $d = 0$ and the result is trivially true. If $\kappa \geq 1$ and $l =0$, again the result is trivially true. If $l \geq 1$ and $u \in K_a$, then reducing $u$ modulo $g_d$ and discarding the last $i_l - i_{l-1}$ terms puts us in the setting with parameters $(\kappa -1, l-1,a)$, where we may use induction. This gives us the surjectivity of $\psi$ and one checks that $J_{\leq -1} \subset \ker \psi$.
\end{proof}

\begin{remark}
    One can prove that, in the setting of Theorem~\ref{regularsequencekernels}, the kernel of $\psi$ is precisely $J_{\leq -1}$ in the cases $l=d$, or with an open condition on the $g_{i_k}$, purely algebraically from the operators defined. Since the argument used in Theorem~\ref{thm:general} is stronger, proving that $\ker \psi = J_{\leq -1}$ when the auxiliary hypersurface
    \[
    Z' = \{g_{i_0} x_0^d + g_{i_1} x_0^{d-i_1} x_1^{i_1} + \dots + g_{i_l} x_1^d \} \subset \p^1_{x_i} \times \Proj S
    \]
    is non-singular, we omit this.
\end{remark}

\begin{proposition}\label{prop:generalWideal}
    Suppose that $l=d$ and $J,U$ are as in Theorem~\ref{regularsequencekernels}. Let $U' = S[W_1,\dots,W_d]$ be the subalgebra of $U$ generated by $W_1,\dots,W_d$. Let $J' \lhd U'$ be the homogeneous ideal generated by the equations
    \[
    g_{k}(Y_k) (W_{k'+1} W_{k''} - W_{k'} W_{k''+1}) - g_{k'}(Y_k) (W_{k+1} W_{k''} - W_k W_{k''+1}) + g_{k''}(Y_k) (W_{k+1} W_{k'} - W_k W_{k'+1})
    \]
    for $0 \leq k < k' < k'' \leq d$. Then the natural map $\phi:U'_{\leq -1} \rightarrow (U/J)_{\leq -1} $ is surjective and has kernel $J'$.
\end{proposition}

\begin{proof}
    The proof is identical to that of Proposition~\ref{prop:universalWideal}, with $\bar U$ being defined as an algebra over $\bar S = S/(g_d(Y_j))$ for the induction step in $d$.
\end{proof}

\begin{thebibliography}{99}
\bibitem{AL} M. Artebani and A. Laface, {\em Hypersurfaces in Mori dream spaces}, J. Algebra 371 (2012), 26--37.
\bibitem{CoxBook} I. Arzhantsev, U. Derenthal, J. Hausen, and A. Laface, {\em Cox rings}, Cambridge Studies in Advanced Mathematics 144, CUP, 2015.

\bibitem{constantinetal} C. Brodie, A. Constantin, J. Gray, A. Lukas, and F. Ruehle, {\em Recent Developments in Line Bundle Cohomology and Applications to String Phenomenology}, in: Nankai Symposium on Mathematical
Dialogues In celebration of S.S.Chern’s 110th anniversary, 2021.
\bibitem{constantin} A. Constantin, {\em Generating Functions for Line Bundle Cohomology Dimensions on Complex Projective Varieties}, Exp. Math. (2024), \href{https://doi.org/10.1080/10586458.2024.2380794}{published online}.
%arXiv:2401.14463
\bibitem{MDS} Y. Hu, S. Keel, {\em Mori dream spaces and GIT}, Michigan Math. J. 48 (2000), 331--348.
\bibitem{hausen} J. Hausen, {\em Cox rings and combinatorics II}, Mosc. Math. J. 8 (2008), 711--757.
\bibitem{lafaceetal} C. Herrera, A. Laface and L. Ugaglia, {\em The Cox ring of an embedded variety}, \href{https://arxiv.org/abs/2411.17370}{arXiv:2411.17370}.
\bibitem{ottem} J. C. Ottem, {\em Birational geometry of hypersurfaces in products of projective spaces}, Math. Z. 280 (2015), 135--148.
\bibitem{pollock} A. Pollock, University of Vienna PhD dissertation, in preparation. 

\end{thebibliography}

\end{document}
